\section{Morphismes universellement ouverts}
\label{section:IV.14-fr}

Les \S\S~14 et 15 sont consacrés à l'étude de la notion de morphisme
\emph{universellement ouvert} (2.4.2). On a déjà vu (2.4.6) qu'un morphisme
\emph{plat} et localement de présentation finie est universellement ouvert, la
réciproque étant inexacte. Au \S~14, on examine d'abord les propriétés des
dimensions des fibres d'un morphisme universellement ouvert $f:X\to Y$~;
lorsque $X$ et $Y$ sont localement noethériens, $f$ se comporte à cet égard
(14.2.1) comme un morphisme plat (cf. 6.1.2), et est en particulier
\emph{équidimensionnel} lorsqu'il est localement de type fini et dominant et que
$X$ est irréductible (14.2.2). Inversement, un morphisme équidimensionnel
localement de type fini $f:X\to Y$ est universellement ouvert lorsqu'on suppose
en outre que $Y$ est \emph{géométriquement unibranche}, et en particulier lorsque
$Y$ est \emph{normal} (critère de Chevalley, 14.4.4). Nous montrons aussi que les
morphismes universellement ouverts localement de type fini $f:X\to Y$ (lorsque
$X$ et $Y$ sont
\oldpage[IV-3]{200}
localement noethériens et $Y$ irréductible) admettent «~suffisamment~» de
\emph{quasi-sections}, i.e. au voisinage d'un point fermé $x$ d'une fibre
$f^{-1}(y)$, il y a un sous-préschéma fermé $X'$ de $X$ contenant $x$ et tel que
la restriction $X'\to Y$ de $f$ soit un morphisme \emph{quasi-fini} (donc à fibres
\emph{discrètes}) et dominant (14.5.3).

Au \S~15 on étudie diverses propriétés des fibres des morphismes universellement
ouverts $f:X\to Y$, localement de type fini, notamment lorsque $X$ et $Y$ sont
localement noethériens. On obtient ainsi en particulier un critère pour qu'un
point $x\in X$ n'appartienne qu'à \emph{une seule} composante irréductible de $X$,
en termes de propriétés de \emph{la fibre} $f^{-1}(f(x))$ de $x$~: il suffit que
$Y$ soit géométriquement unibranche (par exemple normal) au point $f(x)$ et que
$f^{-1}(f(x))$ soit géométriquement ponctuellement intègre au point $x$ (15.3.3)~;
si de plus $Y$ est localement intègre au point $f(x)$, $f$ est \emph{plat} au point
$x$ et $X$ localement intègre au point $x$. On étudie aussi la variation du
\emph{nombre géométrique de composantes connexes} d'une fibre $f^{-1}(y)$~; par
exemple, si $f$ est universellement ouvert et \emph{propre}, et les fibres de $f$
géométriquement réduites, ce nombre est \emph{localement constant} (15.5.7).
Enfin, lorsque $f$ admet une \emph{section} $g$ (ce qui sera le cas lorsque $X$ est
un $Y$-\emph{schéma en groupes}) et que pour tout $y\in Y$, on désigne par $X_y^0$
la composante connexe de la fibre $X_y=f^{-1}(y)$ au point $g(y)$ («~composante
neutre~» dans le cas des groupes), on étudie la réunion $X^0$ des $X_y^0$ pour
$y\in Y$, et on montre (15.6.4) que si $f$ est universellement ouvert et les fibres
$X_y$ géométriquement réduites, alors $X^0$ est un ensemble \emph{ouvert} dans $X$.

\subsection{Morphismes ouverts}
\label{subsection:IV.14.1-fr}

\begin{env}[14.1.1]
\label{IV.14.1.1-fr}
Rappelons (1.10.2) qu'une application continue $\psi:X\to Y$ est dite
\emph{ouverte en un point} $x\in X$ si l'image par $\psi$ de \emph{tout} voisinage
de $x$ dans $X$ est un voisinage de $\psi(x)$ dans $Y$.

On notera que cela \emph{n'implique pas} qu'il existe un système fondamental de
voisinages de $x$ dont les images soient \emph{ouvertes} dans $Y$.
\end{env}

\begin{proposition}[14.1.2]
\label{IV.14.1.2-fr}
Soient $X$, $Y$ deux espaces topologiques, $\psi:X\to Y$ une application
continue, $x$ un point de $X$, $y=\psi(x)$.
\begin{enumerate}[label=(\roman*)]
  \item Si $\psi$ est ouverte en $x$, alors pour toute partie $Y'$ de $Y$ contenant
  $\psi(x)$ la restriction $\psi^{-1}(Y')\to Y'$ de $\psi$ à $Y'$ est ouverte au
  point $x$.
  \item Supposons que $Y$ soit réunion d'une famille localement finie de parties
  fermées $(Y_i)$ et que pour tout $i$ tel que $\psi(x)\in Y_i$, la restriction
  $\psi^{-1}(Y_i)\to Y_i$ de $\psi$ soit ouverte au point $x$~; alors $\psi$ est
  ouverte au point $x$.
  \item Soient $\gamma:X'\to X$ une application continue, $x'$ un point de $X'$~;
  si l'application composée $\psi\circ\gamma:X'\to Y$ est ouverte au point $x'$,
  $\psi$ est ouverte au point $\gamma(x')$.
\end{enumerate}
\end{proposition}

\begin{proof}
Si $U$ est un voisinage de $x$, on a
$\psi(U\cap\psi^{-1}(Y'))=\psi(U)\cap Y'$~; d'où aussitôt (i). Pour prouver (ii),
notons qu'il y a un voisinage $W$ de $\psi(x)$ dans $Y$ ne rencontrant que les
parties fermées $Y_i$ en nombre \emph{fini} qui \emph{contiennent} $\psi(x)$~; donc
$W$ est réunion des $W\cap Y_i$ pour ces indices. Or, si $U$ est un voisinage de
$x$ tel que $\psi(U)\cap Y_i$ soit un voisinage de $\psi(x)$ dans $Y_i$, il existe
un voisinage $V\subset W$ de $\psi(x)$ dans $Y$ tel que pour tous les $i$ tels que
$\psi(x)\in Y_i$,
\oldpage[IV-3]{201}
on ait $V\cap Y_i\subset\psi(U)\cap Y_i$ et comme la réunion des $V\cap Y_i$ pour
ces indices est $V$, on a $V\subset\psi(U)$, donc $\psi(U)$ est un voisinage de
$\psi(x)$ dans $Y$. L'assertion (iii) est triviale.
\end{proof}

\begin{remarks}[14.1.3]
\label{IV.14.1.3-fr}
\begin{enumerate}[label=(\roman*)]
  \item L'ensemble $Z$ des points $x\in X$ où un morphisme est ouvert
  \emph{n'est pas nécessairement ouvert}. Par exemple, soient $K$ un corps, $A$
  l'anneau de polynômes $K[S,T]$, $V$ le plan affine $\operatorname{Spec}(A)$,
  $X$ le sous-préschéma fermé de $V$ «~réunion de la droite $X_1$ définie par
  $T=0$ et de la droite $X_2$ définie par $S=0$~», c'est-à-dire
  $\operatorname{Spec}(A/\mathfrak a)$, où $\mathfrak a=AST$~; prenons
  $Y=X_2=\operatorname{Spec}(A/AS)$ et pour $f$ la projection correspondant à
  l'injection canonique $K[T]\to A/\mathfrak a$~; alors on a $Z=X_2$, qui n'est
  pas ouvert dans $X$.

  \item Soient $X$, $Y$ deux préschémas noethériens irréductibles de points
  génériques $\xi$, $\eta$ respectivement, $f:X\to Y$ un morphisme
  \emph{dominant localement de type fini}~; alors $f$ \emph{est ouvert au point
  $\xi$}. En effet, on peut évidemment se borner au cas où $X$ et $Y$ sont réduits
  (donc intègres) (1.5.4) et affines~; en vertu du théorème de platitude générique
  (6.9.1), il existe un ouvert non vide $V$ de $Y$ tel que la restriction
  $f^{-1}(V)\to V$ de $f$ soit un morphisme \emph{plat}, et on conclut de (2.4.6)
  que cette restriction est un morphisme ouvert. Toutefois, comme il y a des
  morphismes dominants de type fini $f:X\to Y$ (où $X$ et $Y$ sont
  \emph{irréductibles}) qui ne sont pas ouverts (voir par exemple (\textbf{II},
  8.1)), l'ensemble des points où un tel morphisme est ouvert n'est pas
  nécessairement \emph{fermé} dans $X$.

  \item Nous ignorons si, lorsque $X$ et $Y$ sont localement noethériens et
  $f:X\to Y$ un morphisme localement de type fini, l'ensemble des points de $X$
  où $f$ est ouvert est ou non \emph{localement constructible}.
\end{enumerate}
\end{remarks}

\begin{proposition}[14.1.4]
\label{IV.14.1.4-fr}
Soient $X$, $Y$ deux espaces topologiques, $\psi:X\to Y$ une application
continue. Pour tout $y\in Y$, l'ensemble des $x\in\psi^{-1}(y)$ où $\psi$ est
ouvert est une partie fermée de $\psi^{-1}(y)$.
\end{proposition}

\begin{proof}
En effet, supposons que $\psi$ ne soit pas ouvert en un point
$x\in\psi^{-1}(y)$~; il existe alors un voisinage ouvert $V$ de $x$ dans $X$ tel
que $\psi(V)$ ne soit pas un voisinage de $y$~; il en résulte que pour tout
$x'\in V\cap\psi^{-1}(y)$, $\psi$ n'est pas ouvert au point $x'$.
\end{proof}

\begin{remark}[14.1.5]
\label{IV.14.1.5-fr}
Même si $X$ et $Y$ sont localement noethériens et $f$ un morphisme \emph{fini},
il peut se faire que $f$ soit ouvert en tous les points d'une fibre $f^{-1}(y)$,
sans qu'il existe un \emph{voisinage} de $f^{-1}(y)$ en tous les points duquel $f$
soit ouvert. Soient par exemple $K$ un corps, $A$ l'anneau de polynômes
$K[T_1,T_2,T_3]$, $V=\operatorname{Spec}(A)$ l'espace affine à 3 dimensions sur
$K$, $X=\operatorname{Spec}(A/\mathfrak a)$, où $\mathfrak a=\mathfrak b\mathfrak c$,
avec $\mathfrak b=(T_3)$ et $\mathfrak c=(T_1)+(T_2-T_3)$ dans $A$, de sorte que
$X$ est réunion du plan $X_1=\operatorname{Spec}(A/\mathfrak b)$ («~plan
d'équation $T_3=0$~») et de la droite $X_2=\operatorname{Spec}(A/\mathfrak c)$
(«~droite d'équations $T_1=0$, $T_2=T_3$~») qui sont ses composantes
irréductibles. Prenons $Y=X_1$ et soit $f:X\to Y$ la projection correspondant à
l'injection canonique $K[T_1,T_2]\to A/\mathfrak a$~; si $y$ est le point commun
à $X_1$ et $X_2$, $f^{-1}(y)$ est réduit à $y$ et $f$ est ouvert en ce point mais
n'est ouvert en aucun point de $X_2$ dans un voisinage de $y$ et distinct de $y$.
\end{remark}

\begin{env}[14.1.6]
\label{IV.14.1.6-fr}
Dans ce qui suit, le rôle essentiel sera joué par le critère (1.10.3)
caractérisant les morphismes \emph{localement de présentation finie} $f:X\to Y$
qui sont ouverts en un point $x$ par le «~relèvement des générisations~»~: pour
toute générisation $y'$ de $y=f(x)$, il existe $x'\in X$, générisation de $x$, tel
que $y'=f(x')$.
\end{env}

\oldpage[IV-3]{202}
\subsection{Morphismes ouverts et formule des dimensions}
\label{subsection:IV.14.2-fr}

\begin{theorem}[14.2.1]
\label{IV.14.2.1-fr}
Soient $X$, $Y$ deux préschémas localement noethériens, $f:X\to Y$ un morphisme,
$x$ un point de $X$, $y=f(x)$. On suppose que $f$ est ouvert aux points génériques
des composantes irréductibles de $f^{-1}(y)$ contenant $x$. Alors on a la relation
\begin{equation}
\tag{14.2.1.1}
\label{IV.14.2.1.1-fr}
  \dim(\mathcal O_x)=\dim(\mathcal O_y)+
  \dim(\mathcal O_x\otimes_{\mathcal O_y}k(y)).
\end{equation}
\end{theorem}

\begin{proof}
Soit $y'$ une générisation quelconque de $y$ distincte de $y$, et considérons le
sous-préschéma fermé réduit $Y'$ de $Y$ ayant pour espace sous-jacent
$\overline{\{y'\}}$~; alors aucune composante irréductible $X'$ de $f^{-1}(Y')$
contenant $x$ ne peut être contenue dans $f^{-1}(y)$~: en effet, si $x'$ est le
point générique de $X'$, on aurait $f(x')=y\ne y'$, et comme $x'$ est sa seule
générisation dans $f^{-1}(Y')$, cela contredirait l'hypothèse que $f$ est ouvert
au point $x'$, en vertu du fait que $a)$ implique $c)$ dans (1.10.3). On peut donc
appliquer (6.1.2) à l'homomorphisme local d'anneaux noethériens
$\mathcal O_y\to\mathcal O_x$, d'où la conclusion.
\end{proof}

\begin{corollary}[14.2.2]
\label{IV.14.2.2-fr}
Soient $Y$ un préschéma localement noethérien, $f:X\to Y$ un morphisme localement
de type fini, $x$ un point de $X$, $y=f(x)$. Supposons $f$ ouvert aux points
génériques des composantes irréductibles de $f^{-1}(y)$ contenant $x$. Alors $f$
est équidimensionnel au point $x$ dans chacun des deux cas suivants~:
\begin{enumerate}[label=(\roman*)]
  \item $X$ est irréductible et $f$ est dominant.
  \item Les anneaux $\mathcal O_x$ et $\mathcal O_y$ sont équidimensionnels.
\end{enumerate}
\end{corollary}

\begin{proof}
Pour (i), cela résulte de (14.2.1) et de (13.2.3). Pour (ii), cela résulte de
(14.2.1) et de (13.3.6).
\end{proof}

\begin{proposition}[14.2.3]
\label{IV.14.2.3-fr}
Soient $Y$ un préschéma irréductible localement noethérien, $\eta$ son point
générique, $f:X\to Y$ un morphisme localement de type fini, $y$ un point de
$f(X)$, $Z$ une composante irréductible de $f^{-1}(y)$ telle que $f$ soit ouvert
au point générique de $Z$. Alors $Z$ est contenue dans une composante irréductible
$X'$ de $X$ dominant $Y$ et telle que
\[
  \dim(Z)=\dim(X'\cap f^{-1}(y))=\dim(X'\cap f^{-1}(\eta)).
\]
\end{proposition}

\begin{proof}
Cela résulte en effet de (14.2.1) appliqué au point générique de $Z$, et de
(13.2.7).
\end{proof}

\begin{corollary}[14.2.4]
\label{IV.14.2.4-fr}
Avec les notations de (14.2.3), supposons que $f$ soit ouvert en tous les points
de $f^{-1}(y)$. Pour tout $z\in Y$, soient $E(z)$ l'ensemble des dimensions des
composantes irréductibles de $f^{-1}(z)$ et
$d(z)=\sup E(z)=\dim(f^{-1}(z))$. On a alors $E(y)\subset E(\eta)$, d'où
$d(y)\leq d(\eta)$.
\end{corollary}

\begin{proof}
Cela résulte aussitôt de (14.2.3).
\end{proof}

\begin{corollary}[14.2.5]
\label{IV.14.2.5-fr}
Soient $Y$ un préschéma localement noethérien, $f:X\to Y$ un morphisme propre,
$y\in Y$ un point tel que $f$ soit ouvert en tous les points de $f^{-1}(y)$. Alors
la fonction $z\mapsto\dim(f^{-1}(z))$ est constante dans un voisinage de $y$.
\end{corollary}

\begin{proof}
Soient $Y_i$ ($1\leq i\leq n$) les sous-préschémas fermés réduits de $Y$ ayant
pour espaces sous-jacents les composantes irréductibles de $Y$ contenant $y$~;
si $f_i:f^{-1}(Y_i)\to Y_i$ est la restriction de $f$, on sait que chacun des
$f_i$ est propre (\textbf{II}, 5.4.5) et ouvert en tous les points de $f^{-1}(y)$
(14.1.2). Comme la réunion des $Y_i$ est un voisinage de $y$ dans $Y$,
\oldpage[IV-3]{203}
on voit qu'on peut se borner à démontrer le corollaire lorsque $Y$ est
irréductible~; soit $\eta$ son point générique. Il résulte de (14.2.4) que
$d(y)\leq d(\eta)$~; d'autre part, comme $f$ est propre, on déduit de (13.1.5) que
$z\mapsto\dim(f^{-1}(z))$ est semi-continue supérieurement~; en particulier
$d(y)\geq d(\eta)$, donc $d(y)=d(\eta)$. De plus, il y a un voisinage $V$ de $y$
tel que $d(z)\leq d(y)$ pour tout $z\in V$~; mais comme $z$ est spécialisation de
$\eta$, on a d'autre part $d(z)\geq d(\eta)$, d'où finalement $d(z)=d(y)$ pour
$z\in V$.
\end{proof}

\begin{corollary}[14.2.6]
\label{IV.14.2.6-fr}
Soient $Y$ un préschéma localement noethérien, $f:X\to Y$ un morphisme ouvert de
type fini, $y$ un point de $Y$. Soient $y_i$ ($1\leq i\leq n$) les points
génériques des composantes irréductibles de $Y$ contenant $y$. Alors, avec les
notations de (14.2.4)~:
\begin{enumerate}[label=(\roman*)]
  \item Si pour $1\leq i\leq n$, on a $E(y)=E(y_i)$ (resp. $d(y)=d(y_i)$), la
  fonction $E$ (resp. $d$) est constante dans un voisinage de $y$.
  \item Il existe un voisinage ouvert $U$ de $f^{-1}(y)$ tel que la fonction
  $z\mapsto\dim(U\cap f^{-1}(z))$ soit constante dans un voisinage de $y$.
\end{enumerate}
\end{corollary}

\begin{proof}
\begin{enumerate}[label=(\roman*)]
  \item Le même raisonnement que dans (14.2.5) montre qu'on peut se borner à
  considérer le cas où $Y$ est irréductible, de point générique $\eta$. Si $y'$
  est une générisation quelconque de $y$, on a alors (en appliquant (14.2.4) à la
  restriction $f^{-1}(Y')\to Y'$ de $f$, où $Y'=\overline{\{y'\}}$, et utilisant
  (14.1.2)) $E(y)\subset E(y')\subset E(\eta)$ et
  $d(y)\leq d(y')\leq d(\eta)$~; cela montre que la relation $E(y)=E(\eta)$
  (resp. $d(y)=d(\eta)$) entraîne $E(y)=E(y')$ (resp. $d(y)=d(y')$) pour toute
  générisation $y'$ de $y$. Or, en vertu de (9.5.5), les fonctions $E$ et $d$ sont
  localement constructibles, donc il résulte de ($\mathbf{0}_{\mathrm{III}}$,
  9.2.5) appliqué à l'ensemble des $z$ tels que $E(z)=E(y)$ (resp.
  $d(z)=d(y)$) que cet ensemble est un voisinage de $y$.

  \item Pour chacun des $y_i$, $f^{-1}(y_i)$ est un espace noethérien puisque $f$
  est de type fini~; soient $x_{ij}$ ($1\leq j\leq m_i$) les points génériques de
  ses composantes irréductibles, et posons $X'_{ij}=\overline{\{x_{ij}\}}$~;
  montrons que le complémentaire $U$ dans $X$ de la réunion des $X'_{ij}$ qui ne
  rencontrent pas $f^{-1}(y)$ (qui est évidemment un voisinage ouvert de
  $f^{-1}(y)$) répond à la question. En effet, la restriction $f|U:U\to Y$ de $f$
  est un morphisme ouvert de type fini (I, 6.3.5)~; pour tout couple $(i,j)$ tel
  que $x_{ij}\in U$, $x_{ij}$ est point maximal de $U\cap f^{-1}(y_i)$, et il
  résulte de (13.1.1) appliqué à la restriction
  $X'_{ij}\to Y_i=\overline{\{y_i\}}$ de $f$ (compte tenu de (I, 5.2.2) et de
  (1.5.4)) que toutes les composantes irréductibles de
  $X'_{ij}\cap f^{-1}(y)$ ont une dimension
  $\geq\dim(X'_{ij}\cap f^{-1}(y_i))$. Compte tenu de (14.2.3) appliqué à la
  restriction $f^{-1}(Y_i)\to Y_i$ de $f$, qui est un morphisme ouvert (14.1.2),
  $f^{-1}(y)$ est, pour chaque $i$, la réunion des composantes irréductibles de
  $X'_{ij}\cap f^{-1}(y)$ ($1\leq j\leq m_i$), donc
  $\dim(f^{-1}(y))\geq\dim(U\cap f^{-1}(y_i))$~; mais comme $f|U$ est ouvert, on a
  aussi $\dim(f^{-1}(y))\leq\dim(f^{-1}(y_i)\cap U)$ en vertu de (14.2.4)~; on voit
  donc qu'on peut appliquer (i) à $f|U$, d'où la conclusion.
\end{enumerate}
\end{proof}

\begin{remark}[14.2.7]
\label{IV.14.2.7-fr}
L'exemple (13.2.12, (iii)) montre que sous les hypothèses de (14.2.5) ou
(14.2.6, (ii)), on ne peut, dans la conclusion, remplacer la fonction
$z\mapsto\dim(f^{-1}(z))$ par la fonction $z\mapsto E(z)$~; en effet, dans cet
exemple, $f$ est propre et plat (donc universellement ouvert (2.4.6)) et tout
voisinage de l'unique point de $f^{-1}(y_0)$ contient les points génériques des
composantes irréductibles de $f^{-1}(\eta)$.
\end{remark}

\oldpage[IV-3]{204}
\subsection{Morphismes universellement ouverts}
\label{subsection:IV.14.3-fr}

\begin{env}[14.3.1]
\label{IV.14.3.1-fr}
Rappelons (2.4.2) que dire qu'un morphisme $f:X\to Y$ est universellement ouvert
signifie que pour tout morphisme $g:Y'\to Y$, $f_{(Y')}:X_{(Y')}\to Y'$ est
ouvert~; il suffit d'ailleurs qu'il en soit ainsi lorsque
$Y'=Y[T_1,\ldots,T_e]$, pour tout $e$ (8.10.2) (et si $Y$ est localement
noethérien, il suffit donc qu'il en soit ainsi pour tout $Y'$ localement
noethérien).
\end{env}

\begin{proposition}[14.3.2]
\label{IV.14.3.2-fr}
Soit $f:X\to Y$ un morphisme de préschémas. Si $f$ est universellement ouvert,
alors, pour tout morphisme $g:Y'\to Y$, où $Y'$ est irréductible, l'image par
$f'=f_{(Y')}:X_{(Y')}\to Y'$ de toute composante irréductible de $X_{(Y')}$ est
dense dans $Y'$. Réciproquement, si cette condition est vérifiée pour tout $Y'$
irréductible et tout morphisme de type fini $g:Y'\to Y$, et si de plus $f$ est
localement de présentation finie, alors $f$ est universellement ouvert.
\end{proposition}

\begin{proof}
En effet, il résulte de (1.10.4) que si $f$ est universellement ouvert, il vérifie
la condition de l'énoncé. Inversement, supposons que $f$ soit localement de
présentation finie, et montrons que pour tout entier $e$, si l'on pose
$Y''=Y[T_1,\ldots,T_e]$, $f_{(Y'')}$ est ouvert. En effet, soit $Y'$ un
sous-préschéma fermé de $Y''$ ayant pour espace sous-jacent une partie fermée
irréductible de $Y''$~; le morphisme composé $Y'\to Y''\to Y$ est de type fini,
donc toute composante irréductible de $X_{(Y')}$ domine $Y'$ par hypothèse~; on
déduit donc de (1.10.4) que $f_{(Y'')}$ est un morphisme ouvert.
\end{proof}

Cette proposition montre que la définition (\textbf{III}, 4.3.9) coïncide dans
le cas considéré avec la définition générale des morphismes universellement
ouverts donnée dans (2.4.2).

A la notion de morphisme ouvert en un point (14.1.1) correspond de même la
suivante~:

\begin{definition}[14.3.3]
\label{IV.14.3.3-fr}
Soient $f:X\to Y$ un morphisme de préschémas, $x$ un point de $X$. On dit que $f$
est \emph{universellement ouvert au point $x$} si, pour tout morphisme $g:Y'\to Y$,
en posant $X'=X\times_Y Y'$, le morphisme $f'=f_{(Y')}:X'\to Y'$ est ouvert en
tout point $x'$ de $X'$ dont la projection dans $X$ est $x$.
\end{definition}

\begin{remarks}[14.3.3.1]
\label{IV.14.3.3.1-fr}
\begin{enumerate}[label=(\roman*)]
  \item Le raisonnement de (8.10.1) montre (avec les mêmes notations) que si $x$
  est un point de $X$, $x_\lambda$ sa projection dans $X_\lambda$, alors, si
  $f_\mu$ est ouvert au point $x_\mu$ pour tout $\mu\geq\lambda$, $f$ est ouvert
  au point $x$~; il suffit de se borner aux ouverts $U$ de $X$ contenant $x$ et de
  remarquer que l'hypothèse implique que $f_\mu(U_\mu)$ est un voisinage de
  $f_\mu(x_\mu)$, donc $f(v_\mu^{-1}(U_\mu))$ est un voisinage de $f(x)$. On en
  déduit que l'énoncé (8.10.2) est encore exact quand on y remplace
  «~universellement ouvert~» par «~universellement ouvert au point $x$~», et
  «~morphisme ouvert~» par «~morphisme ouvert en tout point $x_n$ de $X_n$ dont
  la projection dans $X$ est $x$~»~: il suffit dans la démonstration de se limiter
  aux ouverts $V$ contenant un $x_n$.

  \item Le résultat de (14.1.4) est encore valable pour un morphisme $f:X\to Y$,
  en y remplaçant «~ouvert~» par «~universellement ouvert~». En effet, supposons
  que $f$ ne soit pas universellement ouvert en un point $x\in f^{-1}(y)$~; il y a
  par suite un morphisme $Y'\to Y$ et un point $x'\in X'$ se projetant en $x$, tels
  que $f'$ ne soit pas ouvert au point $x'$.
\end{enumerate}
\end{remarks}

\oldpage[IV-3]{205}

Or, si $y'=f'(x')$, la projection $f'^{-1}(y')\to f^{-1}(y)$ est un morphisme
ouvert (2.4.10), et il y a, en vertu de (14.1.4), un voisinage $V'$ de $x'$ dans
$f'^{-1}(y')$ où $f'$ n'est pas ouvert~; donc $f$ n'est pas universellement ouvert
aux points de l'image de $V'$ dans $f^{-1}(y)$, qui est un voisinage de $x$ dans
$f^{-1}(y)$.

\begin{proposition}[14.3.4]
\label{IV.14.3.4-fr}
\begin{enumerate}[label=(\roman*)]
  \item Soient $f:X\to Y$, $g:Y\to Z$ deux morphismes, $x$ un point de $X$,
  $y=f(x)$. Si $f$ est universellement ouvert au point $x$ et $g$ universellement
  ouvert au point $y$, alors $g\circ f$ est universellement ouvert au point $x$.
  Inversement, si $g\circ f$ est universellement ouvert au point $x$, $g$ est
  universellement ouvert au point $y$.

  \item Si $f:X\to Y$ est un $S$-morphisme universellement ouvert au point
  $x\in X$, alors, pour tout changement de base $S'\to S$,
  $f_{(S')}:X_{(S')}\to Y_{(S')}$ est universellement ouvert en tout point de
  $X_{(S')}$ au-dessus de $x$.

  \item Pour que $f:X\to Y$ soit universellement ouvert au point $x\in X$, il
  faut et il suffit que $f_{\mathrm{red}}$ le soit.

  \item Soient $f:X\to Y$ un morphisme localement de présentation finie, $x$ un
  point de $X$, $y=f(x)$~; posons $Y_1=\operatorname{Spec}(\mathcal O_y)$,
  $X_1=X\times_Y Y_1$, $f_1=f_{(Y_1)}$. Pour que $f$ soit universellement ouvert
  au point $x$, il faut et il suffit que $f_1$ le soit (on rappelle (I, 3.6.5) que
  $X_1$ est canoniquement identifié à un sous-espace de $X$).
\end{enumerate}
\end{proposition}

\begin{proof}
En effet (ii) est une conséquence évidente de la définition (14.3.3)~; il résulte
aussi de la définition que pour démontrer l'assertion (i), il suffit de le faire
lorsqu'on y supprime partout le mot «~universellement~», et cela résulte alors de
(14.1.2). L'assertion (iii) résulte de (ii) et de ce que le morphisme canonique
$X_{\mathrm{red}}\to X$ est surjectif. Enfin, la condition (iv) est trivialement
nécessaire. D'autre part, si elle est vérifiée, et si $g:Y'\to Y$ est un morphisme
quelconque, $X'=X_{(Y')}$, $f'=f_{(Y')}$, $x'$ un point de $X'$ au-dessus de $x$,
alors $f'$ est localement de présentation finie, et pour voir qu'il est ouvert au
point $x'$, il suffit d'appliquer le critère (1.10.3, $b$). Soient
$y'=f'(x')$, $Y'_1=\operatorname{Spec}(\mathcal O_{y'})$,
$X'_1=X'\times_{Y'}Y'_1$, $f'_1=f_{(Y'_1)}$. Comme $g(y')=y$, le morphisme
composé $Y'_1\to Y'\to Y$ se factorise en $Y'_1\to Y_1\to Y$ (I, 2.4.4), donc
$X'_1=X_1\times_{Y_1}Y'_1$ et $f'_1=(f_1)_{(Y'_1)}$~; la conclusion résulte alors
aussitôt de l'hypothèse que $f'_1$ est ouvert au point $x'$ et de (1.10.3).
\end{proof}

\begin{proposition}[14.3.5]
\label{IV.14.3.5-fr}
Soient $X$, $Y$ deux préschémas, $f:X\to Y$ un morphisme, $x$ un point de $X$.
Soit $(Y_i)_{1\leq i\leq n}$ une famille localement finie de sous-préschémas
fermés de $Y$ telle que l'espace $Y$ soit réunion des $Y_i$, et supposons que pour
tout $i$ tel que $f(x)\in Y_i$, la restriction $f^{-1}(Y_i)\to Y_i$ de $f$ soit un
morphisme universellement ouvert au point $x$~; alors $f$ est universellement
ouvert au point $x$.
\end{proposition}

\begin{proof}
Compte tenu de la définition, cela résulte de la proposition analogue (14.1.2,
(ii)) pour les morphismes ouverts en un point.
\end{proof}

\begin{proposition}[14.3.6]
\label{IV.14.3.6-fr}
Soient $Y$ un préschéma localement noethérien, $f:X\to Y$ un morphisme localement
de type fini. Pour que $f$ soit universellement ouvert en un point $x\in X$, il
faut et il suffit que la condition suivante soit vérifiée~: pour tout morphisme
$g:Y'\to Y$, où $Y'=\operatorname{Spec}(A)$ est le spectre d'un anneau de
valuation discrète, tel que l'image $g(y')$ du point fermé $y'$ de $Y'$ soit égale
à $y=f(x)$, et pour tout point $x'\in X'=X\times_Y Y'$ dont les projections sur
$X$ et $Y'$ sont $x$ et $y'$ respectivement, il existe une générisation $z'$ de
$x'$ dans $X'$ dont la projection dans $Y'$ est le point générique de $Y'$
(autrement dit, il existe une composante irréductible de $X'$ contenant $x'$ et
dominant $Y'$).

\oldpage[IV-3]{206}

De plus, on peut, dans la condition précédente, se borner au cas où $A$ est
complet, a un corps résiduel algébriquement clos, et où $x'$ est rationnel sur
$k(y')$.
\end{proposition}

\begin{proof}
Si $Y'$ est comme dans l'énoncé, la nécessité de la condition résulte de ce que
$f'$ doit être ouvert au point $x'$, et du critère (1.10.3). Pour voir que la
condition est suffisante, considérons un morphisme de type fini $g:Y''\to U$, et
soient $X''=X\times_Y Y''$, $f''=f_{(Y'')}:X''\to Y''$, et $x''$ un point de $X''$
au-dessus de $x$. Posons $y''=f''(x'')$, et soit $t$ une générisation de $y''$
dans $Y''$, distincte de $y''$. Puisque $Y''$ est localement noethérien, il
résulte de (\textbf{II}, 7.1.9) et ($\mathbf{0}_{\mathrm{III}}$, 10.3.1) qu'il
existe un schéma $Y'=\operatorname{Spec}(A)$, où $A$ est un anneau de valuation
discrète complet, dont le corps résiduel est une clôture algébrique de $k(x'')$,
et un morphisme $g:Y'\to Y''$ tels que, si $s$ et $y'$ sont le point générique et
le point fermé de $Y'$, on ait $g(s)=t$ et $g(y')=y''$. Il y a alors un point $x'$
de $X'=X\times_Y Y'=X''\times_{Y''}Y'$ dont les projections dans $X''$ et $Y'$
soient $x''$ et $y'$ respectivement, et qui soit rationnel sur $k(y')$ (I, 3.4.9).
L'hypothèse entraîne qu'il y a une générisation $z'$ de $x'$ dans $X'$ dont la
projection dans $Y'$ soit $s$~; si $z''$ est la projection de $z'$ dans $X''$,
$z''$ est une générisation de $x''$ et sa projection dans $Y''$ est $t$~; on
conclut donc de (1.10.3) que $f''$ est ouvert au point $x''$, donc que $f$ est
universellement ouvert au point $x$ (14.3.3.1, (i)).
\end{proof}

\begin{corollary}[14.3.7]
\label{IV.14.3.7-fr}
Les notations étant celles de (14.3.6)~:
\begin{enumerate}[label=(\roman*)]
  \item Étant donné un point $y\in Y$, pour que $f$ soit universellement ouvert en
  tous les points de $f^{-1}(y)$, il faut et il suffit que pour tout morphisme
  $g:Y'\to Y$, où $Y'$ est le spectre d'un anneau de valuation discrète, et où
  l'image $g(y')$ du point fermé $y'$ de $Y'$ est $y$, toute composante
  irréductible de $X'=X\times_Y Y'$ domine $Y'$.
  \item Pour que $f$ soit universellement ouvert, il faut et il suffit que pour
  tout morphisme $g:Y'\to Y$, où $Y'$ est le spectre d'un anneau de valuation
  discrète, toute composante irréductible de $X'=X\times_Y Y'$ domine $Y'$.
\end{enumerate}
\end{corollary}

\begin{proof}
Il est clair qu'il suffit de prouver (i)~; la nécessité de (i) résulte de (14.3.2)
et sa suffisance de (14.3.6).
\end{proof}

\begin{proposition}[14.3.8]
\label{IV.14.3.8-fr}
Soient $Y$ un préschéma localement noethérien, irréductible, régulier et de
dimension 1 (par exemple le spectre d'un anneau de Dedekind), $f:X\to Y$ un
morphisme localement de type fini, $y$ un point de $Y$. Les conditions suivantes
sont équivalentes~:
\begin{enumerate}[label=\emph{\alph*)}]
  \item $f_{\mathrm{red}}$ est plat en tout point de $f^{-1}(y)$.
  \item $f$ est universellement ouvert dans un voisinage de $f^{-1}(y)$.
  \item $f$ est ouvert dans un voisinage de $f^{-1}(y)$.
  \item Toute composante irréductible de $X$ rencontrant $f^{-1}(y)$ domine $Y$.
\end{enumerate}
\end{proposition}

\begin{proof}
Comme $f_{\mathrm{red}}$ est localement de type fini (1.3.4), $a)$ entraîne que
$f_{\mathrm{red}}$ est plat dans un voisinage de $f^{-1}(y)$ (11.1.1), et il
suffit d'appliquer (2.4.6) dans un tel voisinage pour voir que $a)$ entraîne $b)$.
L'implication $b)\Rightarrow c)$ est triviale, et l'implication
$c)\Rightarrow d)$ résulte de (1.10.4) appliqué à un voisinage de $f^{-1}(y)$.
Reste à voir que $d)$ entraîne $a)$. On peut évidemment, en vertu de (1.3.4), se
borner au cas où $X$ est réduit. La question étant en outre locale sur $X$ et sur
$Y$, on peut supposer $Y=\operatorname{Spec}(A)$ et
$X=\operatorname{Spec}(B)$
\oldpage[IV-3]{207}
affines~; si $X_i$ ($1\leq i\leq n$) sont les sous-préschémas fermés (intègres)
de $X$ ayant pour espaces sous-jacents les composantes irréductibles de $X$,
alors, pour tout $x\in X$, $\mathcal O_{X,x}$, étant réduit, est un sous-anneau du
composé direct des $\mathcal O_{X_i,x}$~; si $y=f(x)$, il suffira de montrer que
chacun des $\mathcal O_{X_i,x}$ est un $\mathcal O_y$-module \emph{sans torsion},
car il en sera alors de même de $\mathcal O_{X,x}$~; comme par hypothèse
$\mathcal O_y$ est un anneau local régulier de dimension 1, c'est-à-dire un anneau
de valuation discrète (\textbf{II}, 7.1.6), il résultera alors de
($\mathbf{0}_{\mathrm I}$, 6.3.4) que $\mathcal O_{X,x}$ est un
$\mathcal O_y$-module \emph{plat}. Mais si $X_i=\operatorname{Spec}(B_i)$, où
$B_i$ est un anneau intègre, l'hypothèse $d)$ entraîne que l'homomorphisme
$A\to B_i$ est \emph{injectif} (I, 1.2.7)~; donc $B_i$ est un $A$-module sans
torsion, et \emph{a fortiori} $\mathcal O_{X_i,x}$ est un $\mathcal O_y$-module
sans torsion.
\end{proof}

\begin{remarks}[14.3.9]
\label{IV.14.3.9-fr}
\begin{enumerate}[label=(\roman*)]
  \item Dans l'énoncé de (14.3.8), on ne peut se dispenser de l'hypothèse que $Y$
  est \emph{régulier}. Avec les notations de (11.7.5), prenons en effet
  $Y=\operatorname{Spec}(A)$, $X=\operatorname{Spec}(\overline A)$, de sorte que
  $f:X\to Y$ est un morphisme \emph{fini surjectif}~; comme $A$ est un anneau
  local intègre de dimension 1, ainsi que $\overline A$, il résulte aussitôt de
  (1.10.4) que le morphisme $f$ est \emph{ouvert}. Cependant $f$ n'est pas
  universellement ouvert (ni \emph{a fortiori} plat), comme le montre (11.7.5).
  On aurait un exemple analogue en prenant pour $Y$ le schéma local au point
  double d'une courbe algébrique ayant un «~point double ordinaire~» et pour $X$
  le normalisé de $Y$.

  \item L'exemple de (14.1.3, (i)) montre que l'ensemble des points de $X$ où un
  morphisme est universellement ouvert n'est pas nécessairement ouvert, le
  morphisme $f$ dans cet exemple étant universellement ouvert en tous les points
  où il est ouvert. L'exemple $f:X\to Y$ vu ci-dessus en (i) montre de même que
  l'ensemble des points où un morphisme est universellement ouvert n'est pas
  nécessairement fermé, car il est immédiat qu'en tous les points de $X$ sauf un
  le morphisme $f$ est universellement ouvert (c'est même un isomorphisme local).
  Il serait intéressant de savoir si l'ensemble des points où un morphisme est
  universellement ouvert est localement constructible.
\end{enumerate}
\end{remarks}

Les deux propositions suivantes nous ont été signalées par M.~Artin~:

\begin{lemma}[14.3.10]
\label{IV.14.3.10-fr}
Soient $A$ un anneau de valuation (non nécessairement discrète), $k$ son corps
résiduel, $K$ son corps des fractions, et posons $S=\operatorname{Spec}(A)$. Soient
$X$ un préschéma irréductible, $f:X\to S$ un morphisme dominant de type fini~; soit
$X_0=X\otimes_A k$ (resp. $X_1=X\otimes_A K$) la fibre de $f$ au point fermé
(resp. au point générique) de $S$. Alors, si $X_0\ne\varnothing$, on a
$\dim(X_0)=\dim(X_1)$.
\end{lemma}

On peut se borner au cas où $X$ est affine, en remplaçant au besoin $X$ par un
ouvert affine contenant un point générique d'une composante irréductible de $X_0$,
de dimension maximale, et utilisant (4.1.1.3). Soit $n=\dim(X_0)\geq 0$~; il résulte
de (13.3.1.1) qu'il existe un voisinage $U$ de $X_0$ dans $X$ et un $S$-morphisme
quasi-fini
\[
  g:U\to S[T_1,\ldots,T_n]=Z
\]
tel que le morphisme restriction
$g_0:U_0=U\cap X_0\to Z_0=\operatorname{Spec}(k[T_1,\ldots,T_n])$ soit fini et
\emph{surjectif}. Par changement de base $\operatorname{Spec}(K)\to S$ et restriction
à l'ouvert $U_1=U\cap X_1$ de $X_1$, on déduit de $g$ un morphisme quasi-fini
$g_1:U_1\to Z_1=\operatorname{Spec}(K[T_1,\ldots,T_n])$. Comme $U_1$ est dense dans
$X_1$, on a $\dim(U_1)=\dim(X_1)$ (4.1.1.3)~; la proposition sera établie, en vertu
de (4.1.2), si nous prouvons que le morphisme $g_1$ est \emph{dominant}. Supposons
le contraire~; il existerait alors un polynôme $F_1\ne 0$ dans
$K[T_1,\ldots,T_n]$ tel que $g_1(U_1)\subset V(F_1)$. Si $\omega$ est une valuation
sur $K$ associée à $A$, et si $(a_\alpha)$ est la famille des coefficients de $F_1$,
on peut, après multiplication de $F_1$ par un élément $\ne 0$ de $K$, supposer que
l'on a $\inf_\alpha(\omega(a_\alpha))=0$~; autrement dit, $F_1$ provient d'un
polynôme $F\in A[T_1,\ldots,T_n]$ (avec lequel il s'identifie) tel que l'image $F_0$
de $F$ dans $k[T_1,\ldots,T_n]$ soit \emph{non nul}. Considérons alors dans $Z$
l'ensemble fermé $V(F)$~; on a $V(F)\cap Z_1=V(F_1)$, et comme $U_1$ est dense
dans $U$ (puisqu'il contient le point générique de $X$), $g(U)\subset V(F)$~; en
particulier, on aurait $g_0(U_0)\subset V(F)\cap Z_0=V(F_0)$~; mais puisque
$F_0\ne 0$, $V(F_0)$ est une partie fermée de $Z_0$ \emph{distincte de $Z_0$} et
l'on aboutit à une contradiction. C.Q.F.D.

\oldpage[IV-3]{208}

\begin{proposition}[14.3.11]
\label{IV.14.3.11-fr}
Soient $f:X\to S$ un morphisme de type fini, $(g_\lambda)_{\lambda\in L}$ une
famille de morphismes universellement ouverts $g_\lambda:Y_\lambda\to S$, et pour
tout $\lambda\in L$, soit $u_\lambda:Y_\lambda\to X$ un $S$-morphisme. Pour tout
$s\in S$, posons $X_s=X\times_S\operatorname{Spec}(k(s))$,
$(Y_\lambda)_s=Y_\lambda\times_S\operatorname{Spec}(k(s))$, et soit
$(u_\lambda)_s:(Y_\lambda)_s\to X_s$ le morphisme $u_\lambda\times 1$ déduit de
$u_\lambda$ par changement de base. Soit $Z(s)$ l'adhérence dans $X_s$ de la
réunion des ensembles $(u_\lambda)_s((Y_\lambda)_s)$, et posons
$d(s)=\dim(Z(s))$. Alors, pour toute générisation $s'$ d'un point $s\in S$, on a
$d(s)\leq d(s')$.
\end{proposition}

\begin{proof}
On sait (\textbf{II}, 7.1.4) qu'il existe un anneau de valuation $A'$ et un
morphisme $h:S'=\operatorname{Spec}(A')\to S$ tels que si $a$ (resp. $b$) est le
point fermé (resp. le point générique) de $S'$, on ait $h(a)=s$, $h(b)=s'$. En
outre, le morphisme projection
$p:X_s\otimes_{k(s)}k(a)\to X_s$ est surjectif et ouvert (2.4.10), donc fait de
$X_s$ un espace quotient de $X_s\otimes_{k(s)}k(a)$ par une relation d'équivalence
ouverte~; pour toute partie $M$ de $X_s$, $p^{-1}(\overline M)$ est donc égal à
l'adhérence $\overline{p^{-1}(M)}$ (Bourbaki, \emph{Top. gén.}, chap.~I\textsuperscript{er},
4\textsuperscript{e} éd., \S~5, n\textsuperscript{o}~3, prop.~7)~; on raisonne de
même pour $X_{s'}$, et compte tenu de (I, 3.4.8), de (4.2.7) et du fait que les
$g_\lambda$ sont universellement ouverts, on voit qu'on peut se ramener à prouver
la proposition sur la situation obtenue après changement de base $S'\to S$.
Supposons donc $S'=S$, $s$ étant le point fermé et $s'$ le point générique de $S$.
L'hypothèse que $g_\lambda$ est ouvert entraîne que toute composante irréductible
de $Y_\lambda$ domine $S$ (1.10.4), donc que son point générique est un point
maximal de $(Y_\lambda)_{s'}$~; si $Z$ désigne l'adhérence dans $X$ de $Z(s')$, on
a donc $u_\lambda(Y_\lambda)\subset Z$, et par suite
$(u_\lambda)_s((Y_\lambda)_s)\subset Z_s=Z\cap X_s$~; autrement dit, on a
$Z(s)\subset Z_s$, d'où $\dim(Z(s))\leq\dim(Z_s)$. Mais en appliquant (14.3.10) à
un sous-préschéma réduit de $X$ ayant pour espace sous-jacent une composante
irréductible de $Z$, on obtient $\dim(Z_s)\leq\dim(Z_{s'})$ (l'inégalité provenant
de ce qu'il peut y avoir des composantes irréductibles de $Z$ ne rencontrant pas
$X_s$). D'autre part, comme $Z(s')$ est par définition fermé dans $X_{s'}$, on a
$Z_{s'}=Z(s')$, ce qui achève la démonstration.
\end{proof}

\begin{remark}[14.3.12]
\label{IV.14.3.12-fr}
Le cas envisagé par M.~Artin était celui où $Y_\lambda=S$ pour tout $\lambda$,
autrement dit le cas où $(u_\lambda)$ est une famille de $S$-sections de $X$. Un
autre cas utile est celui où la famille $(u_\lambda)$ est réduite à un seul
élément~; on peut d'ailleurs toujours se ramener à ce cas en considérant le
préschéma $Y$ somme des $Y_\lambda$ et les morphismes $g:Y\to S$ et $u:Y\to X$
dont les restrictions à chaque $Y_\lambda$ sont respectivement $g_\lambda$ et
$u_\lambda$.
\end{remark}

\begin{proposition}[14.3.13]
\label{IV.14.3.13-fr}
Soient $f:X\to Y$ un morphisme localement de type fini, $y$ un point de $Y$, $x$
un point maximal de la fibre $X_y=X\times_Y\operatorname{Spec}(k(y))$.
Considérons les conditions suivantes~:
\begin{enumerate}[label=\emph{\alph*)}]
  \item $f$ est universellement ouvert au point $x$ (ou encore en tout point de la
  composante irréductible de $X_y$ de point générique $x$ (14.3.3.1, (ii))).
  \item Pour toute composante irréductible $Y_0$ de $Y$ contenant $y$, il existe
  une composante irréductible $Z$ de $X$ contenant $x$, dominant $Y_0$ et telle que
  $\dim_x(X_y)=\dim_x(Z\cap X_y)=\dim(Z\cap X_\eta)$, où $\eta$ est le point
  générique de $Y_0$ (ce qui entraîne que $Z$ est équidimensionnelle sur $Y_0$ au
  point $x$ (13.2.2)).
\end{enumerate}

\oldpage[IV-3]{209}

\begin{enumerate}[label=\emph{\alph*$'$)}]
  \item[b$'$)] Pour tout voisinage ouvert $U$ de $x$ dans $X$ et toute
  générisation $y'$ de $y$, on a
  $\dim(U\cap X_{y'})\geq\dim_x(U\cap X_y)$.
\end{enumerate}
Alors on a les implications $a)\Rightarrow b)\Leftrightarrow b')$.
\end{proposition}

\begin{proof}
Pour montrer que $b)$ implique $b')$, il suffit de remarquer que $y'$ appartient à
une composante irréductible $Y_0$ de $Y$ contenant $y$, de point générique $\eta$~;
prenant $Z$ comme dans $b)$ et notant que le point générique de $Z$ (qui est aussi
celui de $Z\cap X_\eta$ ($\mathbf{0}_{\mathrm I}$, 2.1.8)) est contenu dans $U$, on
a $\dim(U\cap X_{y'})\geq\dim(U\cap Z\cap X_{y'})$, et, en vertu de (13.1.6),
$\dim(U\cap Z\cap X_{y'})\geq\dim(Z\cap X_\eta)$~; d'où l'assertion, puisque
$\dim_x(U\cap X_y)=\dim_x(X_y)$ (4.1.1.3).

Pour prouver que $b')$ implique $b)$, on peut d'abord remplacer $X$ par
$f^{-1}(Y_0)$, donc supposer $Y_0=Y$~; on peut se borner au cas où $X$ et $Y$ sont
affines, puisque $\dim_x(U\cap X_y)=\dim_x(X_y)$ (4.1.1.3). Les composantes
irréductibles $W_i$ du préschéma noethérien $X_\eta$ sont alors en nombre fini, et
le complémentaire de la réunion des $\overline W_i$ (adhérences dans $X$) qui ne
contiennent pas $x$ est un voisinage ouvert $V$ de $x$. En remplaçant $X$ par $V$,
on peut donc supposer que $x\in\overline W_i$ pour tout $i$ (l'hypothèse $b')$
entraîne $\dim(V\cap X_\eta)\geq 0$, donc $x$ appartient à un des $\overline W_i$
pour un $i$ au moins). En outre, les $\overline W_i$ sont exactement les
composantes irréductibles de $X$ qui dominent $Y$ ($\mathbf{0}_{\mathrm I}$,
2.1.8)~; cela étant, si l'on avait, pour chacune de ces composantes $Z$,
$\dim(Z\cap X_\eta)<\dim_x(X_y)$, on en conclurait
$\dim(X_\eta)<\dim_x(X_y)$, contrairement à l'hypothèse $b')$. La relation
$\dim_x(Z\cap X_y)=\dim(Z\cap X_\eta)$ résulte alors de (13.1.6).

Reste à démontrer que $a)$ entraîne $b')$. Tenant compte de (\textbf{II}, 7.1.4)
et de l'invariance des hypothèses et de la conclusion par changement de base (en
vertu de (4.2.7)), on peut se borner au cas où $Y$ est un spectre d'anneau de
valuation, de point fermé $y$ et de point générique $y'$, et où $U=X$.
L'hypothèse que $f$ est ouvert au point $x$ entraîne qu'il existe une composante
irréductible $Z$ de $X$ contenant $x$ et dominant $Y$ (1.10.3). Appliquant
(14.3.10) à un voisinage de $x$ dans $Z$, on en conclut que
$\dim(Z\cap X_{y'})=\dim_x(Z\cap X_y)$~; mais comme $x$ est maximal dans $X_y$,
$Z\cap X_y$ contient la composante irréductible de $X_y$ de point générique $x$,
donc $\dim_x(Z\cap X_y)=\dim_x(X_y)$~; d'autre part, on a
$\dim(X_{y'})\geq\dim(Z\cap X_{y'})$, ce qui achève de prouver $b')$.
\end{proof}

\begin{remark}[14.3.14]
\label{IV.14.3.14-fr}
Nous ignorons si dans (14.3.13), la conclusion subsiste lorsqu'on remplace
l'hypothèse $a)$ par l'hypothèse plus faible que $f$ est ouvert au point $x$. On
peut montrer aisément qu'il suffirait de traiter le cas où $Y$ est spectre d'un
anneau local intègre, dont le point générique est isolé, et où $X$ est un
sous-préschéma fermé du fibré vectoriel $Y[T]$.
\end{remark}

\subsection{Le critère de Chevalley pour les morphismes universellement ouverts}
\label{subsection:IV.14.4-fr}

\begin{theorem}[14.4.1]
\label{IV.14.4.1-fr}
Soient $f:X\to Y$ un morphisme localement de type fini, $y$ un point de $Y$, $x$
un point maximal de la fibre $X_y=X\times_Y\operatorname{Spec}(k(y))$. Supposons
$y$ géométriquement unibranche. Alors les conditions suivantes sont équivalentes~:
\begin{enumerate}[label=\emph{\alph*)}]
  \item $f$ est universellement ouvert au point $x$ (ou encore en tout point de la
  composante irréductible de $X_y$ de point générique $x$ (14.3.3.1, (iii))).

  \oldpage[IV-3]{210}

  \item Si $Y_0$ est l'unique composante irréductible de $Y$ contenant $y$, $\eta$
  son point générique, il existe une composante irréductible $Z$ de $X$ contenant
  $x$, dominant $Y_0$ et telle que
  $\dim_x(Z\cap X_y)=\dim(Z\cap X_\eta)$ (ce qui signifie que $Z$ est
  équidimensionnelle sur $Y_0$ au point $x$ (13.2.2)).
  \item[b$'$)] Pour tout voisinage ouvert $U$ de $x$ dans $X$ et toute
  générisation $y'$ de $y$, on a
  $\dim(U\cap X_{y'})\geq\dim_x(U\cap X_y)$.
\end{enumerate}
Si de plus $Y$ est localement noethérien, ces conditions sont aussi équivalentes à
la suivante~:
\begin{enumerate}[label=\emph{\alph*)}]
  \item[c)] $f$ est ouvert au point $x$.
\end{enumerate}
\end{theorem}

Notons d'abord que puisque $(\mathcal O_y)_{\mathrm{red}}$ est intègre, $y$
n'appartient qu'à une seule composante irréductible $Y_0$ de $Y$. Le fait que $b)$
et $b')$ sont équivalentes et que $a)$ implique $b')$ résulte de (14.3.13)~;
d'autre part, si $Y$ est localement noethérien, on a vu dans (14.2.3) que $c)$
implique $b)$. Il reste donc à montrer que lorsque $y$ est géométriquement
unibranche, $b)$ entraîne $a)$.

\begin{lemma}[14.4.1.1]
\label{IV.14.4.1.1-fr}
Soient $Y$ un préschéma, $Y'=Y[T_1,\ldots,T_n]$, $y'$ un point de $Y'$, $y$ son
image dans $Y$. Si $Y$ est géométriquement unibranche au point $y$, alors $Y'$ est
géométriquement unibranche au point $y'$.
\end{lemma}

En effet, comme pour tout $y\in Y$,
$\operatorname{Spec}(k(y)[T_1,\ldots,T_n])$ est géométriquement régulier sur $k(y)$
($\mathbf{0}_{\mathrm{IV}}$, 17.3.7), le morphisme structural $Y'\to Y$ est lisse
(6.8.1), et il suffit d'appliquer (11.3.14).

\begin{lemma}[14.4.1.2]
\label{IV.14.4.1.2-fr}
Soient $A$ un anneau local intègre unibranche, $B$ un anneau intègre contenant $A$
et entier sur $A$, $\mathfrak n$ un idéal premier de $B$ au-dessus de l'idéal
maximal $\mathfrak m$ de $A$. Alors le morphisme
$\operatorname{Spec}(B_{\mathfrak n})\to\operatorname{Spec}(A)$ est surjectif,
autrement dit, pour tout idéal premier $\mathfrak p$ de $A$, il existe un idéal
premier $\mathfrak q$ de $B$ tel que $\mathfrak q\subset\mathfrak n$ et
$\mathfrak q\cap A=\mathfrak p$.
\end{lemma}

Soient $K$ (resp. $L$) le corps des fractions de $A$ (resp. $B$), $A'$ la clôture
intégrale de $A$, $B'$ le sous-anneau de $L$ engendré par $A'$ et $B$, de sorte
qu'on a un diagramme commutatif d'injections canoniques
\[
  \xymatrix{
    B \ar[r] & B' \\
    A \ar[u] \ar[r] & A' \ar[u]
  }
\]
Comme $B'$ est entier sur $B$, il existe un idéal premier $\mathfrak n'$ de $B'$
tel que $\mathfrak n'\cap B=\mathfrak n$ (Bourbaki, \emph{Alg. comm.}, chap.~V,
\S~2, n\textsuperscript{o}~1, th.~1), et (pour la même raison)
$\operatorname{Spec}(A')\to\operatorname{Spec}(A)$ est surjectif. D'autre part,
comme $A$ est unibranche, $A'$ est un anneau \emph{local}, donc
$\mathfrak n'\cap A'$, qui est au-dessus de l'idéal maximal $\mathfrak m$ de $A$,
est nécessairement égal à l'unique idéal maximal $\mathfrak m'$ de $A'$. En vertu
du second théorème de Cohen--Seidenberg (\emph{loc. cit.}, \S~2,
n\textsuperscript{o}~4, th.~3) le morphisme
$\operatorname{Spec}(B'_{\mathfrak n'})\to\operatorname{Spec}(A')$ est surjectif,
donc il en est de même du composé
$\operatorname{Spec}(B'_{\mathfrak n'})\to\operatorname{Spec}(A')\to
\operatorname{Spec}(A)$~; mais ce morphisme est aussi le composé
$\operatorname{Spec}(B'_{\mathfrak n'})\to\operatorname{Spec}(B_{\mathfrak n})
\to\operatorname{Spec}(A)$, donc le morphisme
$\operatorname{Spec}(B_{\mathfrak n})\to\operatorname{Spec}(A)$ est surjectif.

Ces lemmes étant établis, revenons à la démonstration de l'implication $b)\Rightarrow
a)$ dans (14.4.1). En vertu de (14.3.3.1, (i)), il suffit de prouver que, pour tout
entier $n\geq 0$ et tout point $x'$ de $X'=X[T_1,\ldots,T_n]$ au-dessus de $x$, le
morphisme

\oldpage[IV-3]{211}

$f':X'\to Y'=Y[T_1,\ldots,T_n]$, déduit de $f$ par changement de base, est ouvert
au point $x'$. Compte tenu du lemme (14.4.1.1), de (2.3.4) et (4.2.7), on est donc
ramené à prouver que $f$ est \emph{ouvert} au point $x$~: en outre, il suffit
évidemment (14.1.2, (iii)) de montrer que la restriction de $f$ à un
sous-préschéma fermé de $X$ ayant $Z$ pour espace sous-jacent est ouverte au point
$x$, si bien qu'on peut se borner au cas où $X=Z$ est irréductible. Quitte à
remplacer $X$ par un voisinage ouvert $V$ de $x$ tel que $V\cap X_y$ soit
irréductible, on peut supposer, en vertu de (13.3.1), que le morphisme $f$ se
factorise en
\[
  X\xrightarrow{g}Y''=Y[T_1,\ldots,T_m]\to Y,
\]
où $g$ est quasi-fini, dominant et localement de type fini. Comme le morphisme
structural $Y''\to Y$ est ouvert (2.4.6), on est ramené à prouver que $g:X\to Y''$
est ouvert au point $x$. D'ailleurs, en vertu de (14.4.1.1), $g(x)$ est un point
géométriquement unibranche de $Y''$. On est donc ramené à prouver le lemme suivant~:

\begin{lemma}[14.4.1.3]
\label{IV.14.4.1.3-fr}
Soient $X$, $Y$ deux préschémas irréductibles, $f:X\to Y$ un morphisme localement
quasi-fini et dominant. Si $x\in X$ est tel que $y=f(x)$ soit unibranche sur $Y$,
alors $f$ est ouvert au point $x$.
\end{lemma}

Il suffit de prouver que
$f(\operatorname{Spec}(\mathcal O_{X,x}))=\operatorname{Spec}(\mathcal O_{Y,y})$
(1.10.3). On peut donc se borner, par le changement de base
$\operatorname{Spec}(\mathcal O_{Y,y})\to Y$, au cas où $Y=\operatorname{Spec}(A)$,
où $A$ est un anneau local et $y$ est le point fermé de $Y$ (tenant compte de
(I, 3.6.5) et ($\mathbf{0}_{\mathrm I}$, 2.1.8), qui prouvent que
$X\times_Y\operatorname{Spec}(\mathcal O_{Y,y})$ est irréductible)~; remplaçant $f$
par $f_{\mathrm{red}}$, on peut supposer $X$ et $Y$ réduits, donc intègres.
Remplaçant éventuellement $X$ et $Y$ par des voisinages affines de $x$ et $y$
respectivement, on peut supposer (8.12.9) que le morphisme $f$ se factorise en
$X\xrightarrow{j}X_1\xrightarrow{g}Y$, où $j$ est une immersion ouverte et $g$ un
morphisme \emph{fini} (évidemment dominant)~; comme $X$ et $X_1$ sont affines, $j$
est affine, donc séparé et quasi-compact, et par suite se factorise en
$X\xrightarrow{u}X_2\xrightarrow{h}X_1$, où $X_2$ est l'image fermée de $X$ par
$j$, $h$ l'injection canonique et $u$ une immersion ouverte (I, 9.5.3). En d'autres
termes, on peut supposer que $X_1$ est intègre, ou encore de la forme
$X_1=\operatorname{Spec}(B)$, où $B$ est une $A$-algèbre intègre et finie, contenant
$A$ puisque $g$ est dominant. Si $\mathfrak n$ est l'idéal premier de $B$
correspondant au point $x$, l'hypothèse que $A$ est unibranche implique alors
(14.4.1.2) que le morphisme
$\operatorname{Spec}(B_{\mathfrak n})\to\operatorname{Spec}(A)$ est surjectif,
c'est-à-dire que
$\operatorname{Spec}(\mathcal O_{X,x})\to\operatorname{Spec}(\mathcal O_{Y,y})$
est surjectif. C.Q.F.D.

\begin{corollary}[14.4.2]
\label{IV.14.4.2-fr}
Soient $f:X\to Y$ un morphisme localement de type fini, $y$ un point
géométriquement unibranche de $Y$, $\eta$ le point générique de l'unique composante
irréductible $Y_0$ de $Y$ contenant $y$. Les conditions suivantes sont
équivalentes~:
\begin{enumerate}[label=\emph{\alph*)}]
  \item $f$ est universellement ouvert en tous les points de $X_y$ (ou, ce qui
  revient au même (14.3.3.1, (ii)), aux points maximaux de $X_y$).
  \item Pour tout $x\in X_y$, il existe une composante irréductible $Z$ de $X$
  contenant $x$ et équidimensionnelle sur $Y$ au point $x$ (13.2.2).
  \item[b$'$)] Pour tout $x\in X_y$ et tout voisinage ouvert $U$ de $x$ dans $X$,
  on a $\dim(U\cap X_\eta)\geq\dim_x(U\cap X_y)$.
  \item[b$''$)] Pour tout ouvert $U$ de $X$, on a
  $\dim(U\cap X_\eta)\geq\dim(U\cap X_y)$.
\end{enumerate}

\oldpage[IV-3]{212}

Lorsque de plus $Y$ est localement noethérien, ces conditions sont encore
équivalentes à la suivante~:
\begin{enumerate}[label=\emph{\alph*)}]
  \item[c)] $f$ est ouvert en tous les points de $X_y$ (ou, ce qui revient au même
  (14.3.3.1, (ii)), aux points maximaux de $X_y$).
\end{enumerate}
\end{corollary}

L'équivalence de $a)$ et $c)$ lorsque $Y$ est localement noethérien résulte de
(14.4.1)~; les conditions $b)$ ou $b')$, appliquées aux \emph{points maximaux} de
$X_y$, entraînent $a)$ en vertu également de (14.3.3.1, (ii))~; enfin, $b')$ et
$b'')$ sont équivalentes, car
\[
  \dim(U\cap X_y)=\sup_{x\in X_y}(\dim_x(U\cap X_y)).
\]
Reste à voir que la condition $a)$ entraîne $b)$ et $b')$ en \emph{tout} point
$x\in X_y$. Posons $d=\dim_x(X_y)$, et soit $x'$ le point générique d'une
composante irréductible de $X_y$ contenant $x$ et de dimension $d$. En vertu de
$a)$ et de (14.4.1), il y a une composante irréductible $Z$ de $X$ contenant $x'$
et équidimensionnelle sur $Y$ au point $x'$, donc telle que
$\dim_{x'}(Z\cap X_y)=\dim(Z\cap X_\eta)$. Mais par construction
$\dim_{x'}(Z\cap X_y)=\dim_{x'}(X_y)=d$, et
$\dim_x(Z\cap X_y)\leq\dim(Z\cap X_y)=d$~; compte tenu de (13.1.6), cela prouve
que $Z$ est équidimensionnel sur $Y$ au point $x$~; donc $a)$ entraîne $b)$. De
plus, on a $\dim(X_\eta)\geq\dim(Z\cap X_\eta)=d=\dim_x(X_y)$. Remplaçant $X$ par
un voisinage ouvert $U$ de $x$, on voit ainsi que $a)$ entraîne $b')$. C.Q.F.D.

\begin{corollary}[14.4.3]
\label{IV.14.4.3-fr}
Soient $Y$ un préschéma géométriquement unibranche, $f:X\to Y$ un morphisme
localement de type fini. Les conditions suivantes sont équivalentes~:
\begin{enumerate}[label=\emph{\alph*)}]
  \item $f$ est universellement ouvert.
  \item Pour tout ouvert $U$ de $X$, tout $y\in Y$ et toute générisation $y'$ de
  $y$, on a $\dim(U\cap X_{y'})\geq\dim(U\cap X_y)$.
\end{enumerate}
Si de plus $Y$ est localement noethérien, ces conditions sont aussi équivalentes à
la suivante~:
\begin{enumerate}[label=\emph{\alph*)}]
  \item[c)] $f$ est ouvert.
\end{enumerate}
\end{corollary}

\begin{corollary}[14.4.4]
\label{IV.14.4.4-fr}
(critère de Chevalley). Soit $f:X\to Y$ un morphisme localement de type fini.
\begin{enumerate}[label=(\roman*)]
  \item Si $f$ est équidimensionnel en un point $x\in X$ (13.3.2) et si
  $y=f(x)$ est un point géométriquement unibranche de $Y$, $f$ est
  universellement ouvert au point $x$.
  \item Si $Y$ est géométriquement unibranche, $f$ est universellement ouvert en
  tous les points de $X$ où $f$ est équidimensionnel, et l'ensemble de ces points
  est ouvert dans $X$. En particulier, si $f$ est équidimensionnel, il est
  universellement ouvert.
\end{enumerate}
\end{corollary}

L'assertion (ii) est conséquence triviale de (i), puisqu'on sait déjà que l'ensemble
des points où $f$ est équidimensionnel est ouvert (13.3.2). Quant à l'assertion (i),
elle résulte de ce que l'hypothèse implique que la condition $b)$ de (14.4.1) est
vérifiée au point générique d'une composante irréductible de $X_y$ contenant $x$,
compte tenu de (13.3.1)~; il suffit donc d'appliquer (14.4.1).

\begin{remark}[14.4.5]
\label{IV.14.4.5-fr}
On peut prouver que si $Y$ est localement noethérien, et si toutes les générisations
de $y=f(x)$ sont des points géométriquement unibranches de $Y$ (cf. (6.15.2)),
alors, si $f$ est équidimensionnel au point $x$, il est universellement ouvert
\emph{dans un voisinage de $x$}.
\end{remark}

\oldpage[IV-3]{213}

\begin{corollary}[14.4.6]
\label{IV.14.4.6-fr}
Soient $Y$ un préschéma localement noethérien, $f:X\to Y$ un morphisme localement
de type fini. Soit $y$ un point géométriquement unibranche de $Y$, et supposons en
outre que pour tout $x\in f^{-1}(y)$, l'anneau $\mathcal O_{X,x}$ soit
équidimensionnel. Alors les conditions suivantes sont équivalentes~:
\begin{enumerate}[label=\emph{\alph*)}]
  \item $f$ est équidimensionnel (13.3.2) en tous les points de $f^{-1}(y)$.
  \item $f$ est ouvert en tous les points de $f^{-1}(y)$.
  \item $f$ est universellement ouvert en tous les points de $f^{-1}(y)$.
\end{enumerate}
\end{corollary}

En effet, $c)$ implique trivialement $b)$, et $a)$ implique $c)$ en vertu de
(14.4.4)~; enfin, en raison des hypothèses sur $\mathcal O_y$ et $\mathcal O_x$,
$b)$ implique $a)$ par (14.2.2). Plus généralement~:

\begin{proposition}[14.4.7]
\label{IV.14.4.7-fr}
Soient $Y$ un préschéma localement noethérien, $f:X\to Y$ un morphisme localement
de type fini, $x$ un point de $X$, tel que $y=f(x)$ soit géométriquement unibranche.
Les conditions suivantes sont équivalentes~:
\begin{enumerate}[label=\emph{\alph*)}]
  \item $f$ est équidimensionnel (13.3.2) au point $x$.
  \item L'anneau $\mathcal O_x$ est équidimensionnel et $f$ est ouvert aux points
  génériques des composantes irréductibles de $f^{-1}(y)$ contenant $x$ (donc
  aussi en tout point d'une telle composante).
  \item L'anneau $\mathcal O_x$ est équidimensionnel, et $f$ est universellement
  ouvert aux points génériques des composantes irréductibles de $f^{-1}(y)$
  contenant $x$ (donc aussi en tout point d'une telle composante).
\end{enumerate}
En outre, lorsque ces conditions sont satisfaites, pour tout sous-préschéma fermé
réduit $X_i$ de $X$ ayant pour espace sous-jacent une composante irréductible de $X$
contenant $x$, la restriction $f_i:X_i\to Y$ de $f$ est un morphisme
équidimensionnel au point $x$, et universellement ouvert en tous les points des
composantes irréductibles de $f^{-1}(y)\cap X_i$ qui contiennent $x$.
\end{proposition}

La condition $a)$ implique que $f$ est équidimensionnel en tous les points
génériques des composantes irréductibles de $f^{-1}(y)$ contenant $x$ (13.3.1), et
par suite (14.4.4) universellement ouvert en ces points~; le même raisonnement
appliqué à chaque $f_i$ (compte tenu de (13.3.3)) prouve la dernière assertion de
la proposition, compte tenu de (14.3.3.1, (iii)). En outre, d'après (14.2.1), on a
les relations
\begin{equation}
\tag{14.4.7.1}
\label{IV.14.4.7.1-fr}
  \dim(\mathcal O_{X_i,x})=\dim(\mathcal O_y)+
  \dim(\mathcal O_{X_i,x}\otimes_{\mathcal O_y}k(y))
\end{equation}
\begin{equation}
\tag{14.4.7.2}
\label{IV.14.4.7.2-fr}
  \dim(\mathcal O_{X,x})=\dim(\mathcal O_y)+
  \dim(\mathcal O_{X,x}\otimes_{\mathcal O_y}k(y))
\end{equation}
et puisque $f$ est équidimensionnel au point $x$, il résulte de (13.3.1) que l'on
a $\dim_x(X_i\cap f^{-1}(y))=\dim_x(f^{-1}(y))$, donc (5.2.3)
\[
  \dim(\mathcal O_{X_i,x}\otimes_{\mathcal O_y}k(y))=
  \dim(\mathcal O_{X,x}\otimes_{\mathcal O_y}k(y)).
\]
On conclut donc que $\dim(\mathcal O_{X_i,x})=\dim(\mathcal O_{X,x})$ pour tout
$i$, autrement dit $\mathcal O_{X,x}$ est équidimensionnel, et cela achève de
montrer que $a)$ entraîne $c)$. Il est clair que $c)$ entraîne $b)$~; enfin, $b)$
entraîne la relation (14.4.7.2) en vertu de (14.2.1)~; il résulte alors de (13.3.6)
que $b)$ entraîne $a)$.

\begin{proposition}[14.4.8]
\label{IV.14.4.8-fr}
Soient $Y$ un préschéma noethérien, $f:X\to Y$ un morphisme localement de type
fini, $y$ un point de $Y$, $x$ un point maximal de $X_y=f^{-1}(y)$. Les conditions
suivantes sont équivalentes~:
\begin{enumerate}[label=\emph{\alph*)}]
  \item Le morphisme $f$ est universellement ouvert au point $x$, autrement dit,
  pour tout changement de base $g:Y'\to Y$, on a la propriété $P(Y')$~: pour tout
  point $x'$ de $X'=X\times_Y Y'$ au-dessus de $x$, le morphisme
  $f'=f_{(Y')}:X'\to Y'$ est ouvert au point $x'$.

  \oldpage[IV-3]{214}

  \item[a$'$)] La propriété $P(Y')$ est vraie pour tout morphisme fini
  $g:Y'\to Y$.
  \item[a$''$)] La propriété $P(Y'')$ est vraie pour le normalisé $Y''$ de
  $Y_{\mathrm{red}}$ (\textbf{II}, 6.3.8).
  \item Pour tout point $x''$ de $X''=X\times_Y Y''$ au-dessus de $x$, il existe
  une composante irréductible $Z''$ de $X''$ contenant $x''$ et
  équidimensionnelle sur $Y''$ au point $x''$.
\end{enumerate}
\end{proposition}

Il est trivial que $a)$ implique $a')$. Pour montrer que $a')$ implique $a'')$
notons que l'on peut écrire $Y''=\operatorname{Spec}(\mathcal B)$, où $\mathcal B$
est une $\mathcal O_Y$-Algèbre quasi-cohérente entière sur $\mathcal O_Y$~; comme
$Y$ est noethérien, $\mathcal B$ est limite inductive de ses sous-$\mathcal O_Y$-
Algèbres $\mathcal B_\lambda$ qui sont quasi-cohérentes et de type fini (I, 9.6.6)~;
mais alors les $\mathcal B_\lambda$ sont des $\mathcal O_Y$-Algèbres \emph{finies}
(\textbf{II}, 6.1.2)~; on peut donc écrire
$Y''=\varprojlim Y'_\lambda$, où
$Y'_\lambda=\operatorname{Spec}(\mathcal B_\lambda)$, d'où
$X''=X\times_Y Y''=\varprojlim X'_\lambda$, avec
$X'_\lambda=X\times_Y Y'_\lambda$. En vertu de $a')$, les morphismes
$f'_\lambda:X'_\lambda\to Y'_\lambda$ sont ouverts en tous les points de
$X'_\lambda$ au-dessus de $x$~; on en conclut que $f''$ est ouvert en tous les
points de $X''$ au-dessus de $x$, par (8.10.1) et (14.3.3.1, (i)).

Comme le préschéma $Y''$ est \emph{normal} par définition, le fait que $b)$ entraîne
$a'')$ résulte de (14.4.4) appliqué à la composante irréductible
équidimensionnelle de l'énoncé et à la restriction de $f''$ à cette composante. Il
reste donc à montrer que $a'')$ entraîne $a)$ et $b)$. Compte tenu de (1.10.3), on
peut se borner au cas où $Y=\operatorname{Spec}(\mathcal O_y)$, en notant que le
morphisme canonique $\operatorname{Spec}(\mathcal O_y)\to Y$ est universellement
bicontinu (I, 3.6.5), et d'autre part que
$Y''\times_Y\operatorname{Spec}(\mathcal O_y)$ est le normalisé de
$(\operatorname{Spec}(\mathcal O_y))_{\mathrm{red}}$ comme il résulte de la
permutabilité des opérations de fermeture intégrale et de localisation (Bourbaki,
\emph{Alg. comm.}, chap.~V, \S~1, n\textsuperscript{o}~5, prop.~16). Supposant donc
$Y=\operatorname{Spec}(A)$, où $A$ est un anneau local noethérien, et $y$ le point
fermé de $Y$, on sait ($\mathbf{0}_{\mathrm{IV}}$, 23.2.5) qu'il existe une
factorisation
\[
  Y''\xrightarrow{u}Y_1\xrightarrow{v}Y
\]
du morphisme structural, telle que $v$ soit un morphisme \emph{fini surjectif}, $u$
un morphisme entier, \emph{radiciel et dominant}, donc surjectif (\textbf{II},
6.1.10) et par suite un \emph{homéomorphisme universel} (2.4.5). Si l'on pose
$X_1=X\times_Y Y_1$, la projection $X''\to X_1$ est donc un homéomorphisme, et
l'hypothèse $a'')$ entraîne par suite que $f_1=f_{(Y_1)}:X_1\to Y_1$ est ouvert en
tous les points de $X_1$ au-dessus de $x$. En outre, cela montre que pour prouver
la propriété $b)$, il suffit de prouver la même propriété où l'on remplace $Y''$,
$X''$ et $x''$ par $Y_1$, $X_1$ et un point $x_1$ de $X_1$ au-dessus de $x$~: Mais
$Y_1$ est noethérien et en outre il est géométriquement unibranche puisque $Y''$
est normal et $u$ radiciel (6.15.1)~; la propriété à prouver résulte donc de
(14.4.1). Il reste à montrer que $f$ est universellement ouvert au point $x$, ce
qui résultera du lemme suivant~:

\begin{lemma}[14.4.8.1]
\label{IV.14.4.8.1-fr}
Soit $v:Y_1\to Y$ un morphisme fermé (resp. universellement fermé) et surjectif.
Pour qu'un morphisme $f:X\to Y$ soit ouvert (resp. universellement ouvert) en un
point $x\in X$, il suffit que $f_1=f_{(Y_1)}:X_1=X\times_Y Y_1\to Y_1$ soit ouvert
(resp. universellement ouvert) en tous les points de $X_1$ au-dessus de $x$.
\end{lemma}

La seconde assertion résulte trivialement de la première et du fait que pour tout
changement de base $Y'\to Y$, le morphisme
$v_{(Y')}:Y_1\times_Y Y'\to Y'$ est encore surjectif et est fermé si $v$ est
universellement fermé. Pour prouver la première assertion, considérons

\oldpage[IV-3]{215}

un voisinage ouvert $U$ de $x$ dans $X$~; comme $v$ est fermé et surjectif, pour
que $f(U)$ soit un voisinage de $y=f(x)$, il faut et il suffit que
$v^{-1}(f(U))$ soit un voisinage de $v^{-1}(y)$. Mais si $p:X_1\to X$ est la
projection canonique, on a $v^{-1}(f(U))=f_1(p^{-1}(U))$ (I, 3.4.8), et
l'hypothèse entraîne que $f_1(p^{-1}(U))$ est un voisinage de
$f_1(p^{-1}(x))=v^{-1}(y)$ (I, 3.4.8).

\begin{corollary}[14.4.9]
\label{IV.14.4.9-fr}
Soient $Y$ un préschéma noethérien, $f:X\to Y$ un morphisme localement de type
fini. Les conditions suivantes sont équivalentes~:
\begin{enumerate}[label=\emph{\alph*)}]
  \item $f$ est universellement ouvert, autrement dit, pour tout changement de
  base $Y'\to Y$, le morphisme $f_{(Y')}:X\times_Y Y'\to Y'$ est ouvert.
  \item[a$'$)] Pour tout morphisme fini $Y_1\to Y$, $f_{(Y_1)}$ est ouvert.
  \item[a$''$)] Si $Y''$ est le normalisé de $Y_{\mathrm{red}}$, $f_{(Y'')}$ est
  ouvert.
  \item Pour tout point $x''$ de $X''=X\times_Y Y''$, il existe une composante
  irréductible $Z''$ de $X''$ contenant $x''$ et équidimensionnelle sur $Y''$ au
  point $x''$ (cf. (14.4.10, (ii))).
\end{enumerate}
\end{corollary}

Cela résulte aussitôt de (14.4.8) et (14.1.4).

\begin{remarks}[14.4.10]
\label{IV.14.4.10-fr}
\begin{enumerate}[label=(\roman*)]
  \item L'équivalence des conditions $a)$ et $b)$ dans (14.4.8) (resp. (14.4.9))
  est encore valable pour un préschéma $Y$ quelconque et un morphisme $f$
  localement de type fini. En effet, $a)$ entraîne $b)$ en vertu de (14.4.1)~;
  inversement, $b)$ entraîne que $f''$ est universellement ouvert aux points de
  $X''$ au-dessus de $x$ (resp. en tout point de $X''$) en vertu de (14.4.1), et
  on en conclut la propriété $a)$ en appliquant le lemme (14.4.8.1) au morphisme
  entier et surjectif $Y''\to Y$.

  Il se peut que, dans (14.4.1), pour l'équivalence de $a)$ et $c)$, l'hypothèse
  supplémentaire que $Y$ est noethérien soit superflue (cf. (14.3.14)). S'il en
  est ainsi, les hypothèses noethériennes sont également superflues dans (14.4.2),
  (14.4.3), (14.4.8) et (14.4.9).

  \item On peut donner des exemples de morphismes $f:X\to Y$ ayant les propriétés
  suivantes~: $Y$ est noethérien, régulier et de dimension 2, $f$ est
  \emph{universellement ouvert} et de type fini, $X$ a deux composantes
  irréductibles $X_1$, $X_2$, mais la restriction $X_1\to Y$ de $f$ à l'une
  d'elles \emph{n'est pas un morphisme ouvert}. Le principe de la construction
  s'appuie sur la méthode générale de «~recollement~» qui sera exposée au chap.~V,
  et ne peut donc qu'être esquissé ici. On part d'un point fermé $y$ de $Y$, et on
  considère le $Y$-schéma $Y_1$ obtenu en faisant \emph{éclater} $y$
  (\textbf{II}, 8.1.3)~; si $f_1:Y_1\to Y$ est le morphisme structural, on sait que
  la restriction de $f_1$ à $f_1^{-1}(Y-\{y\})$ est un \emph{isomorphisme} sur
  $Y-\{y\}$ (\emph{loc. cit.}), tandis que la fibre $f_1^{-1}(y)$ est isomorphe à
  $\operatorname{Proj}(S)$, où
  $S=\bigoplus_{k=0}^{\infty}\mathfrak m_y^k/\mathfrak m_y^{k+1}$
  (\textbf{II}, 3.5.3), c'est-à-dire ici à $\mathbf P^1_{k(y)}$~; il résulte de
  (14.4.1) que $f_1$ n'est pas ouvert au point générique de $f_1^{-1}(y)$. D'autre
  part, posons $Y_2=\mathbf P^1_Y$, et soit $f_2:Y_2\to Y$ le morphisme structural~;
  il résulte de (\textbf{II}, 8.4.4) que $f_2$ est plat, donc universellement ouvert
  (2.4.6)~; en outre (\textbf{II}, 3.5.3), $f_2^{-1}(y)$ est isomorphe à
  $\mathbf P^1_{k(y)}$~; il suffit alors de «~recoller~» $Y_1$ et $Y_2$ le long des
  fibres isomorphes $f_1^{-1}(y)$ et $f_2^{-1}(y)$, ce qui donne un morphisme
  $f:X\to Y$ où les composantes irréductibles $X_1$, $X_2$ de $X$ s'identifient
  canoniquement à $Y_1$ et $Y_2$ respectivement, et les restrictions de $f$ à ces
  composantes à $f_1$ et $f_2$.

  Rappelons toutefois (12.1.1.5) que si $Y$ est localement noethérien, $f:X\to Y$
  de type fini et \emph{plat}, alors toute composante irréductible de $X$ est
  \emph{équidimensionnelle} sur $Y$ (et par suite la restriction de $f$ à une telle
  composante est \emph{universellement ouverte} si tous les points de $Y$ sont
  \emph{géométriquement unibranches}).

  \item Rappelons (12.1.2, (i)) qu'il y a des morphismes $f:X\to Y$ ayant les
  propriétés suivantes~: $Y$ est noethérien (non géométriquement unibranche), $f$
  est fini et plat (et même étale (17.6.3)), mais la restriction de $f$ à une
  composante irréductible de $X$ \emph{n'est pas un morphisme ouvert} (bien que $f$
  lui-même le soit d'après (2.4.6)).

  \item Le critère de Chevalley (14.4.4) explique l'importance de la notion de
  morphisme universellement ouvert. Cette notion permet en effet, dans de nombreux
  résultats, plus ou moins classiques, de remplacer une hypothèse de \emph{normalité}
  par l'hypothèse qu'un certain morphisme est universellement ouvert~; l'énoncé plus
  général

  \oldpage[IV-3]{216}

  ainsi obtenu s'appliquera en particulier à des morphismes \emph{plats} (2.4.6),
  dont l'importance en géométrie algébrique va croissant. On peut considérer que
  les énoncés faisant intervenir l'hypothèse qu'un morphisme est universellement
  ouvert sont des généralisations communes d'énoncés faisant intervenir une
  hypothèse de normalité et d'énoncés faisant intervenir une hypothèse de platitude.
\end{enumerate}
\end{remarks}

\subsection{Morphismes universellement ouverts et quasi-sections}
\label{subsection:IV.14.5-fr}

\begin{lemma}[14.5.1]
\label{IV.14.5.1-fr}
Soient $Y$ un préschéma irréductible localement noethérien, $f:X\to Y$ un
morphisme localement de type fini, $x$ un point de $X$, tel que $f$ soit
équidimensionnel (13.3.2) au point $x$~; on pose $y=f(x)$,
$e=\dim_x(f^{-1}(y))$. Soit $X'$ une partie fermée irréductible de $X$ contenant
$x$, $n$ un entier tel que l'on ait $\operatorname{codim}(X',X)\leq n$ et
$\dim_x(X'\cap f^{-1}(y))\leq e-n$. Alors on a nécessairement
$\operatorname{codim}(X',X)=n$, $\dim_x(X'\cap f^{-1}(y))=e-n$, et la restriction $X'\to Y$
de $f$ est un morphisme équidimensionnel en $x$ (et a fortiori dominant).
\end{lemma}

La question étant locale sur $X$, on peut supposer que $f$ est de type fini et
\emph{équidimensionnel}, de sorte que pour \emph{tout} $z\in Y$, \emph{toutes} les
composantes irréductibles de $f^{-1}(z)$ sont de dimension $e$ (13.3.1, $a''$).
Soient $x'$ le point générique de $X'$, $y'=f(x')$,
$Y'=\overline{\{y'\}}=\overline{f(X')}$ et posons $Z=f^{-1}(Y')$. En vertu de
($\mathbf 0$, 14.2.2), on a
\begin{equation}
  \operatorname{codim}(X',Z)\leq n.
  \tag{14.5.1.1}\label{IV.14.5.1.1-fr}
\end{equation}
Et si les deux membres sont égaux, on a nécessairement $\operatorname{codim}(X',X)=n$ et $Z$
contient une composante irréductible de $X$, ce qui entraîne (13.3.1) que
$f(X')$ est dense dans $Y$ et par suite $Z=X$~; donc l'égalité
$\operatorname{codim}(X',Z)=n$ équivaut à la conjonction de l'égalité $\operatorname{codim}(X',X)=n$ et de
la relation $Z=X$.

D'autre part, en raisonnant dans les préschémas réduits de $Y$ et $X$ ayant $Y'$
et $Z$ respectivement pour espaces sous-jacents, on déduit de (5.1.2) et de
(\textbf{I}, 3.6.5) que l'on a
\begin{equation}
  \operatorname{codim}(X'\cap f^{-1}(y'),f^{-1}(y'))=\operatorname{codim}(X',Z).
  \tag{14.5.1.2}\label{IV.14.5.1.2-fr}
\end{equation}
En vertu de l'hypothèse, $f^{-1}(y')$ est biéquidimensionnel (5.2.1) et de
dimension $e$, donc ($\mathbf 0$, 14.3.5), on a, en vertu de (14.5.1.2) et
(14.5.1.1)
\begin{equation}
  e'=\dim(X'\cap f^{-1}(y'))=e-\operatorname{codim}(X',Z)\geq e-n,
  \tag{14.5.1.3}\label{IV.14.5.1.3-fr}
\end{equation}
l'égalité ayant lieu si et seulement si $\operatorname{codim}(X',X)=n$ et si $f(X')$ est dense
dans $Y$. Enfin, (13.1.1), on a $\dim_x(X'\cap f^{-1}(y))\geq e'$, d'où, par
(14.5.1.3),
\begin{equation}
  \dim_x(X'\cap f^{-1}(y))\geq e'\geq e-n.
  \tag{14.5.1.4}\label{IV.14.5.1.4-fr}
\end{equation}
Or, par hypothèse, on a aussi $\dim_x(X'\cap f^{-1}(y))\leq e-n$, d'où les
conclusions de la proposition.

\begin{corollary}[14.5.2]
\label{IV.14.5.2-fr}
Les hypothèses sur $f$, $X$, $Y$, $x$ étant celles de (14.5.1), supposons en
outre que $x$ ne soit pas point maximal de $f^{-1}(y)$. Alors il existe un
voisinage ouvert affine $U$ de $x$ dans $X$, et une section
$g\in\Gamma(U,\mathcal O_X)$ telle que l'ensemble $X'$ des $x'\in U$ tels que
$g(x')=0$ contienne $x$ et ne contienne aucun point maximal de $f^{-1}(y)$. Pour
tout $g$ ayant ces propriétés, $X'$ est équidimensionnel sur $Y$ au point $x$, et
l'on a
\begin{equation}
  \dim_x(X'\cap f^{-1}(y))=e-1 \qquad\text{et}\qquad \operatorname{codim}(X',X)=1.
  \tag{14.5.2.1}\label{IV.14.5.2.1-fr}
\end{equation}
\end{corollary}

\oldpage[IV-3]{217}

On peut se borner au cas où $X=U$ est un voisinage ouvert affine de $x$ tel que
toutes les composantes irréductibles de $f^{-1}(y)$ contiennent $x$. Ces
composantes correspondent aux idéaux premiers minimaux de
$\mathcal O_x/\mathfrak m_y\mathcal O_x$, et par hypothèse ces idéaux sont
distincts de $\mathfrak m_x/\mathfrak m_y\mathcal O_x$ (Bourbaki,
\emph{Alg. comm.}, chap.~II, \S~1, n\textsuperscript{o}~1, prop.~2)~; pour obtenir
un $g\in\Gamma(U,\mathcal O_X)$ vérifiant les conditions de l'énoncé, il suffit de
prendre $g\in\mathfrak j_x$ tel que l'image de $g$ dans $\mathfrak m_x$
n'appartienne à aucun des idéaux premiers précédents. D'ailleurs, on a
$\operatorname{codim}(X',X)\leq 1$ (5.1.8), et comme $X'$ ne contient aucune des composantes
irréductibles de $f^{-1}(y)$ et que celles-ci sont de dimension $e$, on a
($\mathbf 0$, 14.2.2.2) $\dim_x(X'\cap f^{-1}(y))\leq e-1$. Il suffit alors
d'appliquer (14.5.1).

\begin{proposition}[14.5.3]
\label{IV.14.5.3-fr}
Soient $Y$ un préschéma irréductible localement noethérien, $f:X\to Y$ un
morphisme localement de type fini, $x$ un point de $X$. On suppose que $f$ est
équidimensionnel au point $x$ et que $x$ est fermé dans $f^{-1}(f(x))$. Alors il
existe une partie irréductible $X'$ de $X$, localement fermée dans $X$, contenant
$x$ et telle que la restriction $X'\to Y$ de $f$ (où $X'$ est le sous-préschéma
réduit de $X$ ayant $X'$ pour espace sous-jacent) soit un morphisme quasi-fini et
dominant.
\end{proposition}

En effet, avec les notations de (14.5.2), l'hypothèse que $x$ est fermé dans
$f^{-1}(y)$ entraîne que $x$ n'est pas point maximal de $X'\cap f^{-1}(y)$ tant
que $e-1\geq 1$. Il suffit donc d'appliquer (14.5.2) en raisonnant par récurrence
descendante sur $e=\dim_x(f^{-1}(f(x)))$ jusqu'à ce que l'on arrive à $e=1$~;
l'application de (14.5.2) dans ce dernier cas donne un $X'$ tel que
$X'\cap f^{-1}(f(x))$ soit noethérien et de dimension 0, donc fini et discret~;
comme alors $X'$ est équidimensionnel sur $Y$, $X'\cap f^{-1}(f(x'))$ est de
dimension 0 pour tout $x'\in X'$, ce qui entraîne que la restriction $X'\to Y$ de
$f$ est un morphisme quasi-fini (\textbf{II}, 6.2.2).

\begin{corollary}[14.5.4]
\label{IV.14.5.4-fr}
Sous les hypothèses de (14.5.3), supposons en outre que $Y=\operatorname{Spec}(A)$, où $A$ est
un anneau local noethérien intègre et complet, et que $y=f(x)$ soit l'unique point
fermé de $Y$. Alors, il existe un anneau local intègre $A'$, contenant $A$, qui est
une $A$-algèbre finie et possède la propriété suivante~: si l'on pose
$Y'=\operatorname{Spec}(A')$ et $X'=X\times_Y Y'$, il existe une $Y'$-section $h:Y'\to X'$ telle
que le morphisme composé $Y'\xrightarrow{h}X'\to X$ soit une immersion dont
l'image contient $x$.
\end{corollary}

Remplaçant au besoin $X$ par un sous-préschéma réduit irréductible de $X$, on peut,
en vertu de (14.5.3), se borner au cas où le morphisme $f$ est déjà quasi-fini et
dominant. Utilisant (\textbf{II}, 6.2.5), on en déduit que $\mathcal O_x=A'$ est un
anneau intègre qui est une $A$-algèbre finie, et que $X$ est somme disjointe du
sous-préschéma fermé $Y'=\operatorname{Spec}(A')$ et d'un sous-préschéma $Y''$~; le schéma $Y'$
répond à la question, le morphisme composé $Y'\to X'\to X$ n'étant autre que le
morphisme canonique $\operatorname{Spec}(\mathcal O_x)\to X$.

\begin{remark}[14.5.5]
\label{IV.14.5.5-fr}
Si on n'exige pas que dans l'énoncé de (14.5.4), le morphisme $Y'\to X$ soit une
immersion, on peut supposer que $A'$ est en outre \emph{intégralement clos}~: il
suffit en effet de remplacer $A'$ par sa \emph{clôture intégrale} $A_1$, car on
sait ($\mathbf 0$, 23.1.5) que $A_1$ est un $A$-module de type fini.
\end{remark}

\begin{proposition}[14.5.6]
\label{IV.14.5.6-fr}
Soient $Y$ un préschéma localement noethérien, irréductible, régulier et de
dimension 1, $f:X\to Y$ un morphisme localement de type fini, $y$ un point de $Y$.
Les conditions suivantes sont équivalentes~:
\begin{enumerate}[label=\emph{\alph*)}]
  \item $f_{\mathrm{red}}$ est plat en tout point de $f^{-1}(y)$.

  \oldpage[IV-3]{218}

  \item $f$ est universellement ouvert dans un voisinage de $f^{-1}(y)$.
  \item $f$ est ouvert dans un voisinage de $f^{-1}(y)$.
  \item Toute composante irréductible de $X$ rencontrant $f^{-1}(y)$ domine $Y$.
  \item Pour tout point $x\in f^{-1}(y)$, fermé dans $f^{-1}(y)$, et toute
  composante irréductible $X_i$ de $X$ contenant $x$, il existe une partie
  irréductible $X'$ de $X$, localement fermée dans $X$, contenant $x$, contenue
  dans $X_i$, et telle que la restriction $X'\to Y$ de $f$ soit un morphisme
  quasi-fini et dominant.
\end{enumerate}
\end{proposition}

L'équivalence de $a)$, $b)$, $c)$ et $d)$ a déjà été démontrée (14.3.8). Il est
clair que $e)$ entraîne $d)$, car pour toute composante irréductible $X_i$ de $X$
rencontrant $f^{-1}(y)$, il existe dans $X_i\cap f^{-1}(y)$ un point fermé dans cet
espace (5.1.11). Enfin, pour prouver que $c)$ entraîne $e)$, on peut se borner au
cas où $f$ est un morphisme de type fini~; considérons un point fermé $x$ de
$f^{-1}(y)$ et soit $X_i$ une composante irréductible de $X$ contenant $x$. Comme
la restriction $f_i:X_i\to Y$ de $f$ à $X_i$ est un morphisme dominant, il résulte
de l'équivalence de $c)$ et $d)$ pour $f_i$ que ce morphisme est ouvert aux points
génériques de $X_i\cap f^{-1}(y)$. Il résulte donc de (14.2.2) que $f_i$ est
équidimensionnel au point $x$, et on conclut alors à l'aide de (14.5.3).

\begin{remark}[14.5.7]
\label{IV.14.5.7-fr}
Si, dans l'énoncé de (14.5.5), on suppose que $Y=\operatorname{Spec}(A)$, où $A$ est un anneau
de valuation discrète complet, et que $y$ est le point fermé de $Y$, on peut en
outre supposer que $X'=\operatorname{Spec}(A')$, où $A'$ est un anneau de valuation discrète qui
est une $A$-algèbre finie, comme le montre la démonstration de (14.5.4) et le fait
qu'un anneau local intègre régulier et de dimension 1 est un anneau de valuation
discrète (\textbf{II}, 7.1.6).
\end{remark}

\begin{proposition}[14.5.8]
\label{IV.14.5.8-fr}
Soient $Y$ un préschéma localement noethérien, $f:X\to Y$ un morphisme localement
de type fini, $y$ un point de $Y$.

Pour tout $Y$-préschéma $Y'$, posons $X(Y')=\operatorname{Hom}_Y(Y',X)$. Soit $x$ un point de
$f^{-1}(y)$~; afin que $f$ soit universellement ouvert au point $x$, il faut et il
suffit que la condition suivante soit vérifiée~:

Pour tout anneau de valuation discrète complet $A$, de corps résiduel
algébriquement clos, tout morphisme $g:Y'=\operatorname{Spec}(A)\to Y$ tel que l'image par $g$ du
point fermé $y'$ de $Y'$ soit égal à $y$, et tout élément
$u_0\in X(\operatorname{Spec}(k(y')))$ tel que $u_0(y')=x$, il existe un anneau de valuation
discrète $B$, un homomorphisme local $A\to B$ faisant de $B$ une $A$-algèbre finie,
et, si l'on pose $Z'=\operatorname{Spec}(B)$, un élément $u\in X(Z')$ tels que, si $z'$ est le
point fermé de $Z'$, le diagramme
\[
  \xymatrix{
    \operatorname{Spec}(k(z')) \ar[r] \ar[d]^{\ell} & Z' \ar[d]^{u} \\
    \operatorname{Spec}(k(y')) \ar[r]_{u_0} & X
  }
\]
soit commutatif.
\end{proposition}

\oldpage[IV-3]{219}

Notons que si $A$ et $B$ vérifient les conditions de l'énoncé, $B$ est un anneau
de valuation discrète complet (Bourbaki, \emph{Alg. comm.}, chap.~III, \S~3,
n\textsuperscript{o}~3, prop.~7 et chap.~IV, \S~2, n\textsuperscript{o}~2,
cor.~3 de la prop.~9) de corps résiduel isomorphe à celui de $A$, donc
algébriquement clos, et puisque $u(z')=x$, il y a dans $X''=X\times_Y Z'$ un point
$x''$, dont les projections dans $X$ et $Z'$ sont $x$ et $z'$ et qui est rationnel
sur $k(z')$~; en outre, puisqu'il existe une $Z'$-section $v$ de $X''$ telle que
$v(z')=x''$, l'image par $v$ du point générique $s$ de $Z'$ est une générisation
$t$ de $x''$ dont la projection dans $Z'$ est $s$~; en appliquant (14.3.6), on voit
que la condition de l'énoncé est \emph{suffisante}. Prouvons maintenant qu'elle est
\emph{nécessaire}. Posons $X'=X\times_Y Y'$, $f'=f_{(Y')}:X'\to Y'$. Il y a par
hypothèse un point $x'\in X'$ au-dessus de $x$ et de $y'$ et rationnel sur $k(y')$
(\textbf{I}, 3.3.14), donc fermé dans $f'^{-1}(y')$. En vertu de (14.3.6), il y a
une composante irréductible $T'$ de $X'$ contenant $x'$ et dominant $Y'$, donc
((14.3.8) et (14.3.13)) la restriction $T'\to Y'$ de $f'$ est équidimensionnelle au
point $x$. On déduit donc de (14.5.5) qu'il y a une $A$-algèbre finie $B$ qui est
un anneau local intègre et intégralement clos et domine $A$ (donc est un anneau de
valuation discrète) et, en posant $Z'=\operatorname{Spec}(B)$ et
$X''=X\times_Y Z'=X'\times_{Y'}Z'$, une $Z'$-section $v:Z'\to X''$ telle que
l'image du point fermé $z'$ de $Z'$ par le morphisme composé $Z'\to X''\to X'$ soit
$x'$~; le morphisme composé $u:Z'\to X''\to X$ répond donc à la question.

\begin{proposition}[14.5.9]
\label{IV.14.5.9-fr}
Soient $Y$ un préschéma localement noethérien irréductible, $f:X\to Y$ un
morphisme localement de type fini, $y$ un point géométriquement unibranche de $Y$.
Alors les conditions suivantes sont équivalentes~:
\begin{enumerate}[label=\emph{\alph*)}]
  \item $f$ est universellement ouvert en tout point de $f^{-1}(y)$.
  \item $f$ est ouvert en tout point de $f^{-1}(y)$.
  \item Pour toute composante irréductible $Z$ de $f^{-1}(y)$, de point générique
  $z$, il existe une composante irréductible $X_i$ de $X$ contenant $z$ et
  équidimensionnelle sur $Y$ au point $z$.
  \item Pour tout point fermé $x$ de $f^{-1}(y)$, il existe une partie irréductible
  $X'$ de $X$, localement fermée dans $X$, contenant $x$, et telle que la
  restriction $X'\to Y$ de $f$ soit un morphisme quasi-fini et dominant.
\end{enumerate}
\end{proposition}

L'équivalence de $a)$, $b)$ et $c)$ a déjà été démontrée (14.4.2). Pour prouver
que $a)$ entraîne $d)$, notons qu'en vertu de (14.4.2), $a)$ entraîne qu'il existe
une composante irréductible $Z$ de $X$ contenant $x$ et équidimensionnelle sur $Y$
au point $x$~; l'existence de $X'$ provient alors de (14.5.3) appliqué à la
restriction $Z\to Y$ de $f$. Inversement, supposons $d)$ vérifiée~; en vertu du
critère de Chevalley (14.4.4), la restriction $X'\to Y$ de $f$ est un morphisme
universellement ouvert au point $x$, et a fortiori $f$ est ouvert au point $x$~;
$f$ est donc ouvert en tous les points fermés de $X_y=f^{-1}(y)$. Mais $X_y$ est
un $k(y)$-préschéma localement de type fini, donc un préschéma de Jacobson~;
l'ensemble des points fermés de $X_y$ est donc \emph{dense} dans $X_y$ (10.3.1),
et il résulte de (14.1.4) que $f$ est ouvert en \emph{tous} les points de $X_y$, ce
qui achève de prouver que $d)$ entraîne $b)$.

Le résultat suivant a été mis en évidence par D.~Mumford~:

\begin{proposition}[14.5.10]
\label{IV.14.5.10-fr}
Soient $Y$ un préschéma noethérien, $f:X\to Y$ un morphisme universellement ouvert,
surjectif et localement de type fini. Alors il existe un morphisme fini surjectif
$g:Y'\to Y$ tel que, si l'on pose $X'=X\times_Y Y'$ et
$f'=f_{(Y')}:X'\to Y'$, tout point $y'\in Y'$ admette un voisinage ouvert $U'$ tel
qu'il existe une $U'$-section de $f'^{-1}(U')$.
\end{proposition}

Nous démontrerons la proposition en plusieurs étapes.

\oldpage[IV-3]{220}

\emph{I) Réduction au cas où $Y$ est intègre.} -- Si l'on a prouvé la proposition
pour chacun des sous-préschémas réduits $Y_i$ ayant pour espace sous-jacent une
composante irréductible de $Y$, et pour les images réciproques $f^{-1}(Y_i)$, il est
clair que le préschéma $Y'$ \emph{somme} des $Y'_i$ correspondants répondra à la
question. On peut donc supposer $Y$ \emph{intègre} et on va dans ce qui suit se
borner à ce cas. Alors, dans la conclusion, on pourra aussi prendre $Y'$ intègre
(quitte à le remplacer par une composante irréductible convenable).

\emph{II) Caractère local sur $Y$.} -- Nous allons montrer que si l'on peut
recouvrir $Y$ par des ouverts $U_j$ en nombre fini, tels que, pour tout $j$, la
conclusion de la proposition soit vraie pour le morphisme $f^{-1}(U_j)\to U_j$,
restriction de $f$, alors la conclusion est vraie également pour $f$. En effet, on
peut évidemment supposer les $U_j$ affines, de sorte que $U_j=\operatorname{Spec}(A_j)$, où
$A_j$ est un anneau intègre noethérien dont le corps des fractions $K=R(Y)$ est le
corps des fonctions rationnelles sur $Y$. Pour tout $j$, il y a par hypothèse une
$A_j$-algèbre finie intègre $A'_j$ telle que l'homomorphisme $A_j\to A'_j$ soit
injectif (\textbf{I}, 1.2.7) et que le morphisme correspondant
$g_j:U'_j=\operatorname{Spec}(A'_j)\to\operatorname{Spec}(A_j)=U_j$ vérifie les conditions de la proposition
(pour $U_j$ et $f^{-1}(U_j)$). Soit alors $K'$ une extension finie de $K$ contenant
les corps des fractions de tous les $A'_j$ (qui sont des extensions finies de $K$).
Considérons le normalisé $Y''$ de $Y$ dans $K'$ (\textbf{II}, 6.3.8), qui est de la
forme $\operatorname{Spec}(\mathcal B)$, où $\mathcal B$ est une $\mathcal O_Y$-Algèbre entière
quasi-cohérente, fermeture intégrale de $\mathcal O_Y$ dans $K'$ (\textbf{II},
6.3.4). Ces définitions prouvent que pour tout $j$,
\[
  A'_j=\Gamma(U'_j,\mathcal O_{U'_j})
      =\Gamma(U_j,(g_j)_*(\mathcal O_{U'_j}))
\]
s'identifie à une sous-$A_j$-algèbre finie de $\Gamma(U_j,\mathcal B)$~; autrement
dit, $(g_j)_*(\mathcal O_{U'_j})=\mathcal A'_j$ est une
$\mathcal O_{U_j}$-Algèbre cohérente, sous-Algèbre de $\mathcal B|U_j$. Il existe
donc un sous-$\mathcal O_Y$-Module $\mathcal C'_j$ cohérent de $\mathcal B$ tel que
$\mathcal C'_j|U_j=\mathcal A'_j$ (\textbf{I}, 9.4.7). Si l'on pose
$\mathcal C'=\sum_j\mathcal C'_j$, la sous-$\mathcal O_Y$-Algèbre $\mathcal B'$ de
$\mathcal B$ engendrée par $\mathcal C'$ est cohérente puisque $\mathcal B$ est une
$\mathcal O_Y$-Algèbre entière. Montrons alors que $Y'=\operatorname{Spec}(\mathcal B')$ répond à
la question. En effet, il est clair que le morphisme $g:Y'\to Y$ est fini surjectif
et que $Y'$ est intègre~; en outre, pour tout $j$, $\Gamma(U_j,\mathcal B')$ est une
$A'_j$-algèbre finie, autrement dit le morphisme $g^{-1}(U_j)\to U_j$, restriction
de $g$, se factorise en $g^{-1}(U_j)\to U'_j\to U_j$, et comme l'existence locale
de sections est stable par changement de base, cela établit notre assertion, tout
$y\in Y$ appartenant à un $U_j$.

\emph{III) Réduction au cas où $Y$ est intègre, local et géométriquement
unibranche.} -- Supposons d'abord que la proposition ait été prouvée lorsque $Y$
est intègre et \emph{local} (avec $Y'$ intègre), et montrons qu'elle est valable
lorsque $Y$ est intègre (noethérien) affine quelconque. En effet, en vertu de la
réduction II), il suffit de prouver que pour tout point $y\in Y$, la proposition est
vraie pour un voisinage ouvert affine $V$ de $y$ dans $Y$. Soient $Y=\operatorname{Spec}(A)$,
$Y_1=\operatorname{Spec}(A_{\mathfrak p})$, où $\mathfrak p=\mathfrak j_y$, et posons
$X_1=X\times_Y Y_1$~; par hypothèse, il existe un morphisme fini surjectif
$g_1:Y'_1\to Y_1$, où $Y'_1=\operatorname{Spec}(B_1)$, $B_1$ étant une
$A_{\mathfrak p}$-algèbre intègre finie, donc un anneau semi-local, tel que $g_1$
vérifie les conditions de l'énoncé pour $Y_1$ et $X_1$. Si $y'_j$
($1\leq j\leq r$) sont les points fermés de $Y'_1$, il y a donc un recouvrement de
$Y'_1$ par des ouverts $U'_j$ tels que $y'_j\in U'_j$ et qu'il existe une
$U'_j$-section $h'_j$ de $X\times_Y U'_j$ ($1\leq j\leq r$). Le
$A_{\mathfrak p}$-module $B_1$ admet un système fini de générateurs de la forme
$z_i/s$ (avec $s\in A-\mathfrak p$,

\oldpage[IV-3]{221}

$z_i$ entiers sur $A$), que l'on peut supposer (en multipliant au besoin $s$ par un
élément de $A$) être des éléments du corps des fractions de $B_1$, entiers sur
$A_s$, de sorte que si $V$ est l'ouvert affine $D(s)=\operatorname{Spec}(A_s)\subset Y$, $Y'_1$
s'identifie à $Y_1\times_V V'$, où $V'$ est le spectre de la $A_s$-algèbre finie
engendrée par les $z_i/s$~; $g:V'\to V$ est donc un morphisme fini surjectif et
$g_1=g_{(Y_1)}$. De plus, appliquant la méthode de (8.1.2, $a)$), on peut supposer
que chacun des $U'_j$ est l'image réciproque par $p:Y'_1\to V'$ d'un ouvert $W'_j$
de $V'$, tel que les $W'_j$ recouvrent $V'$ (8.3.11), et que chacune des sections
$h'_j$ est de la forme $(v'_j)_{(Y_1)}$, où $v'_j$ est une $W'_j$-section de
$X\times_Y W'_j$ (8.8.2, (i)). On est donc bien ramené à prouver la proposition
lorsque $Y=\operatorname{Spec}(A)$, $A$ étant un anneau local intègre et noethérien. On sait alors
qu'il existe une $A$-algèbre intègre finie $B$, ayant même corps des fractions que
$A$, telle que $A\subset B$ et que $\operatorname{Spec}(B)$ soit géométriquement unibranche
(($\mathbf 0$, 23.2.5) et (6.15.5))~; comme le morphisme
$\operatorname{Spec}(B)\to\operatorname{Spec}(A)$ est surjectif, on peut évidemment remplacer $Y$ par
$\operatorname{Spec}(B)$ et $X$ par $X\otimes_A B$ pour prouver la proposition. Raisonnant comme
au début de la réduction III), on peut donc supposer $A$ local, intègre et
$Y=\operatorname{Spec}(A)$ géométriquement unibranche.

\emph{IV) Réduction au cas où $X$ est intègre, affine, et $f$ quasi-fini,
surjectif, birationnel et universellement ouvert.} -- Supposons donc $Y=\operatorname{Spec}(A)$
intègre, local et géométriquement unibranche. Il existe alors un sous-préschéma
irréductible $X_0$ de $X$ tel que la restriction $f_0:X_0\to Y$ de $f$ soit un
morphisme quasi-fini et dominant et que $f_0(X_0)$ contienne le point fermé $y$ de
$Y$ (14.5.9)~; comme $Y$ est géométriquement unibranche, il résulte de (14.4.1) que
$f_0$ est encore \emph{universellement ouvert}. Comme en outre on peut supposer
$X_0$ réduit, donc intègre, on voit qu'on peut, en remplaçant $X$ par $X_0$,
supposer que $X$ est \emph{intègre} et $f$ \emph{quasi-fini et dominant}, et tel que
$f(X)$ contienne $y$~; comme $f$ est \emph{ouvert} et que tout ouvert de $Y$
contenant $y$ est égal à $Y$, $f$ est surjectif.

Soient $\xi$, $\eta$ les points génériques de $X$ et $Y$ respectivement~;
$k(\xi)$ est donc une extension \emph{finie} de $K=k(\eta)$. Par suite (4.6.8), il
y a une extension \emph{finie} $K'$ de $K$ telle que
$(\operatorname{Spec}(k(\xi)\otimes_K K'))_{\mathrm{red}}$ soit géométriquement réduit sur $K$ et
que ses composantes irréductibles soient géométriquement irréductibles~; il en
résulte ((4.5.9) et (4.6.1)) que les corps résiduels de
$\operatorname{Spec}(k(\xi)\otimes_K K')$ sont des extensions finies, primaires et séparables de
$K'$, donc sont \emph{égaux à} $K'$. Appliquant de nouveau ($\mathbf 0$, 23.2.5) et
(6.15.5), il existe une $A$-algèbre intègre \emph{finie} $B$, ayant $K'$ pour
corps des fractions, contenant $A$ et telle que $\operatorname{Spec}(B)$ soit géométriquement
unibranche~; on peut donc de nouveau remplacer $Y$ par $\operatorname{Spec}(B)$ et $X$ par
$(X\otimes_A B)_{\mathrm{red}}$ pour prouver la proposition~; la réduction III)
permet alors de supposer encore $A$ local, intègre et $Y=\operatorname{Spec}(A)$ géométriquement
unibranche, de point fermé $y$. En outre, $f$ est alors un morphisme \emph{quasi-fini,
surjectif et universellement ouvert}~; chacune des composantes irréductibles $X_i$
de $X$ ($1\leq i\leq m$) domine $Y$ (1.10.4) et est \emph{birationnelle} sur $Y$.
Soit alors $x$ un point de $f^{-1}(y)$ et soit $X_i$ une composante irréductible de
$X$ contenant $x$~; désignons encore par $X_i$ le sous-préschéma réduit (donc
intègre) de $X$ ayant $X_i$ pour espace sous-jacent, et soit $f_i:X_i\to Y$ la
restriction de $f$. Appliquant de nouveau (14.4.1) au morphisme quasi-fini et
dominant $f_i$, on voit que $f_i$ est universellement ouvert~; si $U$ est un ouvert
affine de $X_i$ contenant $x$, $f_i(U)$ est donc

\oldpage[IV-3]{222}

ouvert dans $Y$ et contient $y$, donc est égal à $Y$. On peut ainsi remplacer $X$
par $U$ pour prouver la proposition.

\emph{V) Fin de la démonstration.} -- Nous supposons donc $Y$ local, intègre et
géométriquement unibranche, $X$ intègre et affine, $f$ quasi-fini, surjectif,
birationnel et universellement ouvert~; il en résulte que $f$ est automatiquement
\emph{séparé}. En vertu du Main theorem (8.12.6), il existe une factorisation
$X\xrightarrow{j}Z\xrightarrow{u}Y$, où $j$ est une immersion ouverte et $u$ un
morphisme \emph{fini}. Remplaçant en outre $Z$ par l'image fermée de $j$
(\textbf{I}, 9.5.10), on peut supposer que $Z$ est \emph{intègre}~; comme $u$ est
fini et birationnel et $Y$ géométriquement unibranche, il résulte de
(\textbf{III}, 4.3.5 et 4.3.4) que $u$ (et par suite $f$) est un morphisme
\emph{radiciel}~; comme $f$ est universellement ouvert et surjectif, c'est donc un
\emph{homéomorphisme universel}. Par suite (2.4.5, (ii)) $f$ est un morphisme
\emph{fini}. Mais alors on répond aux conditions de l'énoncé avec $Y'$ \emph{intègre}
en prenant simplement $Y'=X$ et $g=f$, la $Y'$-section étant le morphisme diagonal
$\Delta_f$. C.Q.F.D.

Ce résultat peut être utilisé pour développer des critères de «~descente~» de
diverses propriétés par les morphismes universellement ouverts et surjectifs.
Signalons en particulier le critère suivant, dû à D.~Mumford~:

\begin{corollary}[14.5.11]
\label{IV.14.5.11-fr}
Soient $Y$ un préschéma localement noethérien, $f:X\to Y$ un morphisme localement
de type fini, $g:Y'\to Y$ un morphisme localement de type fini, universellement
ouvert et surjectif~; posons $X'=X\times_Y Y'$, $f'=f_{(Y')}:X'\to Y'$. Alors,
pour que $f$ soit affine, il faut et il suffit que $f'$ le soit.
\end{corollary}

On doit seulement prouver que la condition est suffisante. La question étant locale
sur $Y$, on peut supposer $Y$ affine, donc noethérien. En vertu de (14.5.10), il
existe un morphisme fini surjectif $h:Y_1\to Y$ tel que, si l'on pose
$Y'_1=Y'\times_Y Y_1$, et $g'=g_{(Y_1)}:Y'_1\to Y_1$, tout point de $Y_1$ admette
un voisinage ouvert $U_1$ tel qu'il existe une $U_1$-section de $g'^{-1}(U_1)$. Si
l'on pose $X_1=X\times_Y Y_1$, la projection canonique $p:X_1\to X$ est un
morphisme fini surjectif, donc, si l'on prouve que $X_1$ est un schéma affine, il en
résultera d'abord que $X$ est quasi-compact, donc noethérien, puis que $X$ est affine
en vertu du théorème de Chevalley (\textbf{II}, 6.7.1). Il suffit par suite de
prouver que le morphisme $f_1=f_{(Y_1)}:X_1\to Y_1$ est affine. Or, si l'on pose
\[
  X'_1=X_1\times_{Y_1}Y'_1=X'\times_{Y'}Y'_1
  \qquad\text{et}\qquad
  f'_1=(f_1)_{(Y'_1)}=f'_{(Y'_1)}:X'_1\to Y'_1,
\]
$f'_1$ est affine en vertu de l'hypothèse. On est donc ramené à prouver le corollaire
en y remplaçant $Y$, $X$, $f$ et $Y'$ par $Y_1$, $X_1$, $f_1$ et $Y'_1$,
autrement dit, il suffit de prouver que $f$ est affine lorsqu'on fait de plus, dans
l'énoncé de (14.5.11), l'hypothèse que tout point de $Y$ admet un voisinage ouvert
$U$ tel qu'il existe une $U$-section de $g^{-1}(U)$. La question étant locale sur
$Y$, on peut même supposer qu'il existe une $Y$-section $s$ de $Y'$. Or, on a le
lemme élémentaire suivant (valable dans toute catégorie admettant des produits
fibrés)~:

\begin{lemma}[14.5.11.1]
\label{IV.14.5.11.1-fr}
Soient $f:X\to S$, $g:Y\to S$ deux morphismes, $p_1:X\times_S Y\to X$,
$p_2:X\times_S Y\to Y$ les projections canoniques. Si $s:S\to Y$ est une section
de $g$, alors $s'=(1,s\circ f)_S$ est une section de $p_1$ et $X$, muni des
morphismes $f:X\to S$ et $s':X\to X\times_S Y$, s'identifie au produit des
$Y$-préschémas $S$ et $X\times_S Y$ pour les morphismes $s:S\to Y$ et
$p_2:X\times_S Y\to Y$.
\end{lemma}

\oldpage[IV-3]{223}

\begin{proof}
C'est un cas particulier de (\textbf{I}, 3.3.11), où l'on remplace le diagramme par
\[
  \xymatrix{
    X \ar[r]^{p_1} \ar[d]_{f} & X\times_S Y \ar[r]^{s'} \ar[d]^{p_2} & X \ar[d]^{f} \\
    S \ar[r]_{s} & Y \ar[r]_{g} & S
  }
\]
\end{proof}

Appliquant ce lemme en y remplaçant $S$, $Y$ par $Y$, $Y'$, on voit que l'on peut
écrire $f=(f')_{(Y)}$ pour le changement de base $s:Y\to Y'$, donc $f$ est affine
puisque $f'$ l'est. C.Q.F.D.

Une variante de ce critère est la suivante~:

\begin{corollary}[14.5.12]
\label{IV.14.5.12-fr}
Avec les notations et hypothèses générales de (14.5.11), soit $\mathcal L$ un
$\mathcal O_X$-Module inversible. Supposons qu'il existe une partie fermée $Z$ de
$X$, propre sur $Y$ (\textbf{II}, 5.4.10) et telle que le préschéma induit par $X$
sur l'ouvert $X-Z$ soit normal. Alors, pour que $\mathcal L$ soit ample relativement
à $f$, il faut et il suffit que
$\mathcal L'=\mathcal L\otimes_{\mathcal O_Y}\mathcal O_{Y'}$ soit ample
relativement à $f'$.
\end{corollary}

Gardons les notations de la démonstration de (14.5.11). Posons
$\mathcal L_1=\mathcal L\otimes_{\mathcal O_Y}\mathcal O_{Y_1}$~; il suffira de
prouver que $\mathcal L_1$ est ample relativement à $f\circ p$, en vertu de
(\textbf{III}, 2.6.2)~; mais $f\circ p=h\circ f_1$, et, compte tenu de
(\textbf{II}, 4.6.13, (v)), il suffira de prouver que $\mathcal L_1$ est ample
relativement à $f_1$. La question étant locale sur $Y_1$, on peut de nouveau
supposer que $g'$ admet une section $s$~; si
$\mathcal L'_1=\mathcal L'\otimes_{\mathcal O_{Y'}}\mathcal O_{Y'_1}$, on peut
alors écrire, en vertu du lemme (14.5.11.1),
$\mathcal L_1=\mathcal L'_1\otimes_{\mathcal O_{Y'_1}}\mathcal O_{Y_1}$ pour le
changement de base $s:Y_1\to Y'_1$~; la conclusion résulte donc de deux
applications de (\textbf{II}, 4.6.13, (iii)).

% IV3_B31_P223_END_BEFORE_SECTION_15
